\documentclass[11pt,a4paper]{article}

\usepackage[utf8]{inputenc}
\usepackage[T1]{fontenc}
\usepackage[french]{babel}
\usepackage{amsmath,amssymb,amsthm}
\usepackage[ruled,vlined,linesnumbered]{algorithm2e}
\usepackage[margin=2.4cm]{geometry}
\usepackage{booktabs}
\usepackage{enumitem}
\usepackage{xcolor}
\usepackage{tikz}
\usepackage{pgfplots}
\pgfplotsset{compat=1.18}
\usetikzlibrary{positioning,arrows.meta,shapes.geometric,fit,backgrounds}
\usepackage{listings}
\usepackage[hidelinks]{hyperref}

% ---------------------------------------------------------------- algorithmes
\renewcommand{\algorithmcfname}{Algorithme}
\SetKwInput{KwIn}{Entrées}
\SetKwInput{KwOut}{Sorties}
\SetKwFor{For}{pour}{faire}{fin pour}
\SetKwFor{While}{tant que}{faire}{fin tant que}
\SetKwFor{ForEach}{pour chaque}{faire}{fin pour chaque}
\SetKwIF{If}{ElseIf}{Else}{si}{alors}{sinon si}{sinon}{fin si}
\SetKw{Return}{retourner}
\SetKw{KwTo}{à}
\SetKw{Break}{sortir}
\SetKwComment{Comment}{$\triangleright$\ }{}
\DontPrintSemicolon

% ---------------------------------------------------------------- théorèmes
\theoremstyle{plain}
\newtheorem{theorem}{Théorème}
\newtheorem{proposition}[theorem]{Proposition}
\newtheorem{lemma}[theorem]{Lemme}
\theoremstyle{remark}
\newtheorem{remark}{Remarque}

% ---------------------------------------------------------------- notations
\newcommand{\Zk}{Z_k}
\newcommand{\Nk}{N_k}
\newcommand{\Dk}{D_k}
\newcommand{\Ecal}{\mathcal{E}}
\newcommand{\Scal}{\mathcal{S}}
\newcommand{\Rcal}{\mathcal{R}}
\newcommand{\Acal}{\mathcal{A}}
\newcommand{\ILP}{\textsc{Pline}}
\newcommand{\qs}{q^{\star}}

% ---------------------------------------------------------------- couleurs
\definecolor{ceff}{RGB}{0,102,153}     % coupe d'efficacité
\definecolor{cdom}{RGB}{190,90,20}     % coupe de dominance
\definecolor{cgris}{RGB}{160,160,160}
\definecolor{cfond}{RGB}{235,242,247}
\definecolor{cfonddom}{RGB}{250,238,225}

% ---------------------------------------------------------------- listings
\lstset{
  language=Python, basicstyle=\ttfamily\footnotesize,
  keywordstyle=\color{ceff}\bfseries, commentstyle=\color{cgris}\itshape,
  stringstyle=\color{cdom}, numbers=left, numberstyle=\tiny\color{cgris},
  numbersep=7pt, frame=leftline, framesep=6pt, rulecolor=\color{cgris},
  showstringspaces=false, breaklines=true, columns=fullflexible,
  literate={é}{{\'e}}1 {è}{{\`e}}1 {ê}{{\^e}}1 {à}{{\`a}}1 {û}{{\^u}}1
           {ô}{{\^o}}1 {î}{{\^i}}1 {ç}{{\c c}}1 {É}{{\'E}}1
}

\title{\bfseries La coupe d'efficacité\\[2pt]
\large Une hybridation exacte--heuristique pour l'optimisation d'une fonction
linéaire fractionnaire sur l'ensemble efficace d'un MOILFP}
\author{}
\date{}

\begin{document}
\maketitle
\vspace{-1.6\baselineskip}

\begin{abstract}
\noindent
On considère le problème (P) : maximiser une fonction linéaire fractionnaire
$f$ sur l'ensemble $\Ecal$ des solutions efficaces d'un programme linéaire
fractionnaire multi-objectif en nombres entiers. Ce mémoire présente
\emph{une seule} contribution et la traite complètement : une \textbf{coupe
d'efficacité} qui permet de retrancher du relâché servant à borner $\qs$ la
région dominée par un point \textbf{efficace}, alors que la coupe classique
ne l'autorise qu'autour d'un point \emph{dominé}.

L'intérêt est une hybridation au sens fort. La recherche heuristique produit,
comme sous-produit, une archive de points efficaces certifiés ; le
théorème~\ref{th:effcut} en fait la \emph{source des coupes de la borne
exacte}. Le flux d'information s'inverse : ce n'est plus la partie exacte qui
certifie la partie heuristique, c'est la partie heuristique qui fournit à la
partie exacte ce dont elle manque. La réserve -- les solutions alternatives
de même vecteur critères -- se lève par une condition de clôture qui coûte
\emph{un seul programme de faisabilité} par point d'archive.

Une proposition en découle, sans équivalent pour la coupe de dominance : si
le relâché se \emph{vide}, l'optimalité est prouvée sans qu'aucune borne
n'ait été évaluée. Un lemme réduit ensuite le coût de la clôture à un
\emph{programme linéaire}.

Les mesures sont rapportées par \textbf{paires} -- même instance, même
graine, les deux variantes -- et jamais en comparant deux médianes
indépendantes, qui laissent croire à un gain par instance qui n'existe pas.
Sur le régime cible, l'optimalité prouvée passe de $58/80$ à $78/80$
exécutions, soit $20$ preuves gagnées et \emph{aucune} perdue, avec un coût
en baisse ($-12$ appels au solveur entier par paire en médiane). Sur un lot
figé de $90$ instances, $2$ preuves gagnées, aucune perdue, valeur identique
partout. Un banc à budget \emph{déterministe} -- plafonné en nombre d'appels
et non en secondes -- confirme la reproductibilité sur $90$ instances sur
$90$ et l'absence de tout recul.

La méthode est enfin confrontée aux deux méthodes exactes voisines, dont
aucune ne résout notre problème -- l'une optimise une fonction linéaire sur
l'ensemble efficace d'un MOILFP, l'autre une fonction fractionnaire sur
celui d'un MOILP -- ce qui impose de les affronter chacune sur son propre
terrain. Leurs deux exemples publiés sont reproduits exactement, et celle de
Zerdani \& Moulaï est réimplémentée, tableau du simplexe compris, en
arithmétique rationnelle. Sur $24$ instances de sa classe, elle certifie
$12$ optima quand notre méthode en prouve $24$ ; mais elle dispose d'un
raccourci qui la rend imbattable sur $9$ d'entre elles, et nous le disons.

Les limites sont énoncées avec la même précision que les gains : la méthode
est sans effet au-delà de $n\ge 20$, et à budget d'appels égal le lemme
n'achète aucune preuve supplémentaire -- c'est un gain de coût, non de
capacité. Une coupe disjonctive sans aucune variable binaire, obtenue par
programme générateur, y est construite et vérifiée : elle ne gagne que
$0{,}11$ point d'écart garanti, ce qui \emph{infirme} le diagnostic que nous
avions nous-mêmes posé sur l'origine de ce verrou. Deux illustrations numériques entièrement détaillées, à $2$ puis
$3$ variables, exposent le mécanisme pas à pas.
\end{abstract}

\tableofcontents
\bigskip

% ============================================================================
\section{Introduction}

\subsection{Le problème}

Soit le programme linéaire fractionnaire multi-objectif en nombres entiers
\[
  \text{(MOILFP)}\qquad
  \max\ \Zk(x)=\frac{c_k^{\top}x+a_k}{d_k^{\top}x+b_k},\quad k=1,\dots,p,
  \qquad \text{s.c. } Ax\le b,\ x\in\mathbb{Z}^n_+ ,
\]
et $\Ecal$ son ensemble de solutions efficaces. Étant donnée une seconde
fonction fractionnaire $f=(c^{\top}x+a)/(d^{\top}x+b)$, on cherche
\[
  \text{(P)}\qquad \qs=\max_{x\in\Ecal} f(x).
\]
La difficulté n'est pas $f$, qui se traiterait par Dinkelbach ; elle est le
domaine, $\Ecal$, qui n'est ni convexe, ni connexe, ni décrit par des
contraintes. Toute méthode doit donc \emph{représenter} $\Ecal$ d'une
manière ou d'une autre.

\subsection{Ce qui manque aux méthodes existantes}

Les méthodes exactes du domaine énumèrent implicitement $\Ecal$. Elles sont
correctes, et sans recours dès que le front s'épaissit : interrompues, elles
ne rendent \emph{aucune} garantie. Une méthode qui encadre $\qs$ à tout
instant, elle, reste exploitable à budget fini.

Un tel encadrement demande deux choses. Une borne inférieure : la valeur
$q=f(\bar x)$ d'un point $\bar x$ \emph{certifié efficace} -- c'est le
travail de la recherche. Et une borne supérieure : un ensemble
$\Rcal\supseteq\Ecal$ sur lequel maximiser un substitut linéaire. La qualité
de la borne tient entièrement à la \emph{finesse} de $\Rcal$, donc à ce qu'on
sait en retrancher.

\subsection{Le verrou, en une phrase}

On ne sait retrancher que ce qui est dominé par un point \textbf{dominé}
déjà rencontré. Or les points dominés sont rares là où la borne en aurait le
plus besoin : la profondeur médiane des chaînes de réparation vaut $1$, et
sur les deux instances illustratives de ce mémoire la recherche n'en produit
\emph{aucun}. La coupe classique n'a rien à mordre.

\subsection{Contribution}

\begin{enumerate}[leftmargin=*,itemsep=3pt]
\item \textbf{Théorème~\ref{th:effcut} (coupe d'efficacité).} On peut couper
      autour d'un point \emph{efficace} : la région retranchée ne contient,
      parmi les points efficaces, que ceux de \emph{même vecteur critères}.
\item \textbf{Théorème~\ref{th:archbound} (borne sur le relâché d'archive).}
      Sous une condition de clôture, la borne resserrée reste valide sur ce
      relâché plus petit.
\item \textbf{Proposition~\ref{prop:vide} (relâché vide).} Si le relâché se
      vide, $\qs=q$. La coupe de dominance ne peut jamais produire cette
      preuve, puisqu'elle \emph{préserve} $\Ecal$.
\item \textbf{La clôture coûte un programme de faisabilité}, pas un optimum
      (\S\ref{sec:cloture}), et ses deux issues sont l'une et l'autre utiles.
\item \textbf{Théorème~\ref{th:hauteur} (clôture par la hauteur).} Un
      \emph{majorant} $\le 0$ de la hauteur de la région retranchée établit
      la clôture. Comme il se calcule sur la relaxation continue, un
      programme \emph{linéaire} remplace le programme en nombres entiers.
      C'est ce qui fait passer le coût apparié de $+10$ à $-12$ appels par
      exécution sur le régime cible.
\item \textbf{Un seul vivier, deux sources} (\S\ref{sec:vivier}) : mettre les
      deux coupes en concurrence sous un plafond global, et non leur allouer
      des budgets séparés -- point où une mise en \oe uvre naïve perd
      $28$ preuves sur $90$.
\item \textbf{La coupe d'efficacité comme opérateur de recherche}
      (\S\ref{sec:divers}) : couper autour de toute l'archive rend un point
      \emph{garanti neuf} en vecteur critères, et cet usage-là n'exige
      aucune condition de clôture. Sur $24$ paires à budget déterministe,
      $6$ instances améliorées, $\mathbf{0}$ dégradée.
\item \textbf{Une comparaison sur le terrain de chacun}
      (\S\ref{sec:comparaison}) : réimplémentation de Zerdani \& Moulaï avec
      son simplexe rationnel à tableau, reproduction des deux exemples
      publiés, et match à unité de coût indépendante de la machine.
\item \textbf{Un protocole de mesure qui se contrôle lui-même}
      (\S\ref{sec:determ}) : lecture appariée plutôt que comparaison de
      médianes, budget plafonné en \emph{appels} et non en secondes, et
      vérification du déterminisme par double exécution. Sans ce protocole,
      deux des conclusions de ce mémoire auraient été fausses -- nous le
      montrons plutôt que de l'affirmer.
\end{enumerate}

% ============================================================================
\section{Cadre et briques réutilisées}
\label{sec:cadre}

Cette section fixe les notations et rappelle, sans les redémontrer, les
seuls résultats antérieurs dont la suite a besoin. Elle est délibérément
brève : tout le reste du mémoire est nouveau.

\subsection{Notations et hypothèses}

$\Scal=\{x\in\mathbb{Z}^n_+ : Ax\le b\}$ est le domaine entier,
$\Nk(x)=c_k^{\top}x+a_k$ et $\Dk(x)=d_k^{\top}x+b_k$, $Z(x)$ le vecteur
critères. Un point $y$ \emph{domine} $x$ si $Z(y)\ge Z(x)$ et $Z(y)\ne Z(x)$ ;
$\Ecal$ est l'ensemble des $x\in\Scal$ non dominés. Un rationnel $q$ est
toujours écrit $q=P/Q$ sous forme irréductible avec $Q>0$.

\paragraph{Hypothèses.} (A1) $d_k\ge 0$ et $b_k\ge 1$ pour tout $k$, ainsi
que pour $f$ ; (A2) $A\ge 0$ à colonnes non nulles et $b\ge 0$. (A1) entraîne
$\Dk(x)\ge 1>0$ partout ; (A2) rend $\Scal$ non vide et borné. Toutes les
données sont entières.

\begin{remark}[intégralité]\label{rem:int}
Les quantités $e_k$ et $QN-PD$ introduites plus bas sont \emph{entières}.
Les tests d'arrêt portant sur elles sont donc exacts, sans tolérance
numérique. C'est ce qui rend sûres les équivalences
« $e_k(x)>0 \iff e_k(x)\ge 1$ » et « $f(x)>q \iff QN(x)-PD(x)\ge 1$ »
utilisées de façon répétée.
\end{remark}

\subsection{Trois briques}

\paragraph{Brique 1 -- test d'efficacité intégral.} Pour $\bar x\in\Scal$, en
posant $\bar N_k=\Nk(\bar x)$, $\bar D_k=\Dk(\bar x)$ et
\[
  e^{\bar x}_k(x)\ :=\ \bar D_k\,\Nk(x)-\bar N_k\,\Dk(x)
   \ =\ \bar D_k\,\Dk(x)\,\bigl(\Zk(x)-\Zk(\bar x)\bigr),
\]
la quantité $\theta(\bar x)=\max\bigl\{\sum_k e^{\bar x}_k(x)\ :\
e^{\bar x}_k(x)\ge 0\ \forall k,\ x\in\Scal\bigr\}$ est un entier
$\ge 0$, et $\bar x\in\Ecal\iff\theta(\bar x)=0$. Si $\theta(\bar x)>0$,
tout optimum \emph{domine} $\bar x$.

Deux conséquences servent ici. Le test est \emph{exact}, donc l'archive
produite par la recherche est faite de points \emph{prouvés} efficaces --
c'est la précondition du théorème~\ref{th:effcut}. Et $e^{\bar x}_k$ est
entière, de même signe que $\Zk(x)-\Zk(\bar x)$, d'où
\begin{equation}\label{eq:eint}
  \Zk(x)>\Zk(\bar x)\iff e^{\bar x}_k(x)\ge 1 .
\end{equation}

\paragraph{Brique 2 -- borne resserrée.} Soit $q=P/Q$ atteinte par un point
efficace, $\Rcal$ un relâché contenant $\Ecal$,
$U\ge\max_{x\in\Rcal}\{QN(x)-PD(x)\}$ et
$D^{+}=\min\{D(x):x\in\Rcal,\ QN(x)-PD(x)\ge 0\}$. Alors
\begin{equation}\label{eq:borne}
  \qs\ \le\ q+\frac{U}{Q\,D^{+}},
  \qquad\text{et } U\le 0 \Rightarrow \qs=q .
\end{equation}

\paragraph{Brique 3 -- coupe de dominance.} Si $\bar x$ est
\emph{dominé}, la région $\{x\in\Scal : Z(x)\le Z(\bar x)\}$ ne contient
aucun point efficace, et se retranche par la disjonction
\begin{equation}\label{eq:disj}
  \exists k:\ e^{\bar x}_k(x)\ge 1
  \quad\Longleftrightarrow\quad
  e^{\bar x}_k(x)\ge 1-M_k(1-u_k)\ \ \forall k,\quad
  \textstyle\sum_k u_k\ge 1,\quad u\in\{0,1\}^p .
\end{equation}
C'est la seule opération de resserrement disponible avant ce mémoire. Son
coût est de $p$ variables binaires par coupe -- il faut le garder en tête,
c'est lui qui impose un plafond (\S\ref{sec:vivier}).

% ============================================================================
\section{Le verrou : la coupe de dominance n'a rien à mordre}
\label{sec:verrou}

La précondition de la brique 3 n'est pas négociable : couper autour d'un
point efficace le retirerait de $\Rcal$, et $\Rcal$ ne contiendrait plus
$\Ecal$. Or elle contraint la coupe à ne se nourrir que de ce que les
chaînes de réparation traversent.

\paragraph{Un constat mesuré, pas une intuition.} La campagne mesure une
profondeur médiane de chaîne de réparation égale à $1$ : la recherche ne
franchit qu'un point dominé par mouvement. Dans le régime à $\Ecal$ épais --
celui qui résiste -- l'archive compte des dizaines de points efficaces
certifiés quand les points dominés se comptent sur les doigts. Et sur les
deux instances illustratives de ce mémoire, la recherche en collecte
\textbf{zéro} : la réparation atteint un point efficace du premier coup.

\begin{center}
\begin{tabular}{lcc}
\toprule
& instance $\mathcal{I}_2$ & instance $\mathcal{I}_3$\\
\midrule
$|\Scal|$ & 22 & 36\\
$|\Ecal|$ & 5 & 5\\
points dominés \emph{existants} & 17 & 31\\
points dominés \emph{collectés par la recherche} & $\mathbf{0}$ & $\mathbf{0}$\\
points d'archive disponibles & $\mathbf{5}$ & $\mathbf{4}$\\
\bottomrule
\end{tabular}
\end{center}

\smallskip\noindent
La ressource est là -- l'archive -- mais le théorème interdit de s'en servir.
Deux obstacles à lever : montrer qu'une coupe autour d'un point efficace ne
détruit que peu, et rendre ce « peu » vérifiable.

% ============================================================================
\section{La coupe d'efficacité}
\label{sec:effcut}

\subsection{Le théorème}

\begin{theorem}[Coupe d'efficacité]\label{th:effcut}
Soient $a\in\Ecal$ et $x\in\Ecal$ avec $Z(x)\ne Z(a)$. Alors il existe $k$
tel que $e^a_k(x)\ge 1$. Par conséquent, pour tout $A\subseteq\Ecal$,
\[
  \Ecal\ \subseteq\ \Rcal_A\ \cup\
  \bigl\{x\in\Scal:\ \exists a\in A,\ Z(x)=Z(a)\bigr\},
  \qquad
  \Rcal_A:=\bigl\{x\in\Scal:\ \forall a\in A,\ \exists k,\ e^a_k(x)\ge 1\bigr\}.
\]
\end{theorem}

\begin{proof}
Supposons $\Zk(x)\le \Zk(a)$ pour tout $k$. Comme $Z(x)\ne Z(a)$, le point
$a$ dominerait $x$, ce qui contredit l'efficacité de $x$. Il existe donc $k$
avec $\Zk(x)>\Zk(a)$, c'est-à-dire $e^a_k(x)\ge 1$
par~\eqref{eq:eint}. L'inclusion s'obtient en appliquant ceci à chaque
$a\in A$, pour les $x$ efficaces dont le vecteur critères diffère de tous
ceux de $A$.
\end{proof}

\noindent
La différence avec la brique 3 tient en une ligne, et c'est toute la
contribution :

\begin{center}
\begin{tabular}{lll}
\toprule
& \textbf{coupe de dominance} & \textbf{coupe d'efficacité}\\
\midrule
précondition sur $a$ & $a$ \emph{dominé} & $a$ \emph{efficace}\\
forme algébrique & \multicolumn{2}{c}{identique : la disjonction~\eqref{eq:disj}, $p$ binaires}\\
points efficaces perdus & aucun & ceux de \emph{même vecteur critères}\\
source & chaînes de réparation & \textbf{archive}\\
$\Rcal$ peut-il se vider ? & \textbf{jamais} ($\Rcal\supseteq\Ecal$) & \textbf{oui}\\
\bottomrule
\end{tabular}
\end{center}

\subsection{La borne reste valide}

\begin{theorem}[Borne sur le relâché d'archive]\label{th:archbound}
Soient $A\subseteq\Ecal$ un ensemble de points efficaces \emph{certifiés},
$q=P/Q$ la valeur de l'incumbent, et supposons la \textbf{condition de
clôture}
\begin{equation}\label{eq:cloture}
  \forall a\in A:\qquad
  \max\{f(x)\ :\ x\in\Scal,\ Z(x)=Z(a)\}\ \le\ q .
\end{equation}
Soient $U\ge\max_{x\in\Rcal_A}\{QN(x)-PD(x)\}$ et
$D^{+}=\min\{D(x):x\in\Rcal_A,\ QN(x)-PD(x)\ge 0\}$. Alors
\[
  \boxed{\ \qs\ \le\ q+\frac{U}{Q\,D^{+}}\ }
\]
\end{theorem}

\begin{proof}
Soit $x^{\star}\in\Ecal$ atteignant $\qs$. Par le
théorème~\ref{th:effcut}, deux cas. Ou bien $Z(x^{\star})=Z(a)$ pour un
$a\in A$ : alors $\qs=f(x^{\star})\le q$ par~\eqref{eq:cloture}, et
l'inégalité est acquise puisque $U\ge 0$. Ou bien
$x^{\star}\in\Rcal_A$ : l'argument de~\eqref{eq:borne} s'applique mot pour
mot à $\Rcal_A$, soit
$f(x^{\star})-q=\bigl(QN(x^\star)-PD(x^\star)\bigr)/\bigl(Q D(x^\star)\bigr)
\le U/(Q D^{+})$.
\end{proof}

\subsection{Une preuve inaccessible à la coupe de dominance}

\begin{proposition}[Relâché vide]\label{prop:vide}
Sous les hypothèses du théorème~\ref{th:archbound}, si $\Rcal_A=\emptyset$
alors $\qs=q$.
\end{proposition}

\begin{proof}
Par le théorème~\ref{th:effcut},
$\Ecal\subseteq\Rcal_A\cup\{x:\exists a\in A,\ Z(x)=Z(a)\}$. Si
$\Rcal_A=\emptyset$, tout point efficace a donc le vecteur critères d'un
point de $A$, et la clôture~\eqref{eq:cloture} donne $\qs\le q$. Comme
$q\le\qs$ ($q$ est atteinte par un point efficace), il y a égalité.
\end{proof}

\noindent
Cette proposition n'a \emph{pas} d'équivalent pour la coupe de dominance.
Celle-ci préserve $\Ecal$ par construction : $\Rcal$ contient toujours
$\Ecal$ et ne peut donc jamais se vider tant que $\Ecal$ est non vide. Seule
la coupe d'efficacité peut \emph{épuiser} le relâché, et c'est une preuve
d'optimalité obtenue \textbf{sans qu'aucune borne n'ait été évaluée}.

Le cas se présente exactement là où $\Ecal$ est mince -- le régime à
corrélation élevée entre critères, où la campagne mesure une médiane
$|\Ecal|=1$ -- puisque l'archive y couvre $\Ecal$ en quelques points. Ne pas
le traiter coûte cher : dans une version antérieure de notre mise en \oe
uvre, un relâché vide était lu comme un \emph{échec} de la borne, et
l'optimalité prouvée tombait à $56$ instances sur $90$ au lieu de $86$, la
totalité de l'écart provenant d'instances à front mince.

% ============================================================================
\section{La condition de clôture coûte un programme, pas un optimum}
\label{sec:cloture}

La condition~\eqref{eq:cloture} semble demander, pour chaque $a\in A$, un
maximum de $f$ sur $\{x:Z(x)=Z(a)\}$ -- donc un Dinkelbach complet par point
d'archive, ce qui ruinerait l'intérêt de la coupe. Il n'en est rien :
\emph{on n'a pas besoin du maximum, seulement de savoir s'il dépasse $q$}.
C'est une question de \textbf{faisabilité}, et un seul programme y répond.

\begin{equation}\label{eq:clot}
  \text{(C}_a\text{)}\qquad
  \text{trouver } x\in\Scal
  \quad\text{tel que}\quad
  e^a_k(x)=0\ \ \forall k,
  \qquad Q\,N(x)-P\,D(x)\ \ge\ 1 .
\end{equation}

\noindent
Les deux contraintes sont exactes grâce à la remarque~\ref{rem:int} : les
$e^a_k$ étant entières, $e^a_k(x)=0$ équivaut à $\Zk(x)=\Zk(a)$ ; et
$QN-PD$ étant entière, « $\ge 1$ » équivaut à $f(x)>q$. Deux issues, toutes
deux utiles :

\begin{itemize}[leftmargin=*,itemsep=3pt]
\item \textbf{(C$_a$) infaisable} $\Rightarrow$ la clôture est établie pour
      $a$ : \emph{on pose la coupe} ;
\item \textbf{(C$_a$) faisable} $\Rightarrow$ le point rendu a le même
      vecteur critères qu'un point efficace, il est donc \emph{lui-même
      efficace}, et il bat l'incumbent. On l'empoche : la coupe n'est pas
      posée, mais la borne \emph{inférieure} monte, ce qui est un pas de
      Dinkelbach gratuit.
\end{itemize}

\noindent
Il n'y a donc pas de branche perdante. En pratique la seconde issue ne s'est
jamais produite dans nos mesures : sur $143\,220$ paires de points efficaces
vérifiées par énumération exhaustive, aucune ne partage de vecteur critères.
La clôture est de fait gratuite -- mais elle est \emph{vérifiée} et non
supposée, et c'est ce qui rend la coupe sûre.

\subsection{Un programme linéaire suffit souvent}
\label{sec:lemme}

Le programme (C$_a$) est en nombres entiers. On peut faire moins cher, et
le gain n'est pas marginal : c'est un appel entier de moins par point
d'archive.

\begin{theorem}[Clôture par la hauteur]\label{th:hauteur}
Posons la \emph{hauteur de la région} retranchée par une coupe en $a$ :
\[
  h(a)\ =\ \max\ \bigl\{\,Q\,N(x)-P\,D(x)\ :\ x\in\Scal,\ Z(x)\le Z(a)\,\bigr\}.
\]
Si $h(a)\le 0$, la condition de clôture~\eqref{eq:cloture} est établie
en $a$.
\end{theorem}

\begin{proof}
$\{x : Z(x)=Z(a)\}\subseteq\{x : Z(x)\le Z(a)\}$. Si le maximum de
$QN-PD$ sur le second ensemble est $\le 0$, il l'est sur le premier,
c'est-à-dire $f(x)\le q$ pour tout $x$ de même vecteur critères que $a$.
\end{proof}

\noindent
L'intérêt tient au fait que seul un \emph{majorant} de $h(a)$ est requis :
un majorant $\le 0$ suffit à conclure. On le calcule sur la
\textbf{relaxation continue} de $\Scal$, qui la contient -- donc un
\emph{programme linéaire} remplace un programme en nombres entiers. Le
maximum est de plus ramené au plancher entier, $QN-PD$ étant entière.

\begin{remark}[ce que ce lemme est, et ce qu'il n'est pas]\label{rem:lemme}
Il ne change \emph{aucune} coupe : même vivier, même ordre, mêmes coupes
posées. Seul le moyen d'autoriser la coupe change. Le nombre d'appels au
solveur entier ne peut donc que baisser -- par construction, et pas
seulement en moyenne. C'est un gain de \textbf{coût}. Ce n'est pas un gain
de \textbf{capacité} : la section~\ref{sec:res} mesure qu'à budget d'appels
\emph{égal}, il n'achète aucune preuve supplémentaire.
\end{remark}

\noindent
Le comportement mesuré est nettement bimodal : là où le lemme opère il
établit presque toutes les clôtures ($24$ sur $24$, $22$ sur $22$, $21$ sur
$21$), là où il n'opère pas il n'en établit aucune. Il n'y a pas de régime
intermédiaire.

% ============================================================================
\section{Un seul vivier, deux sources}
\label{sec:vivier}

Les deux coupes ont exactement la même forme -- la disjonction~\eqref{eq:disj}
-- donc exactement le même coût, $p$ binaires chacune. Leur allouer des
budgets séparés serait doublement fautif.

D'abord parce qu'un plafond existe : au-delà d'une quarantaine de coupes,
chaque coupe supplémentaire alourdit le programme entier plus qu'elle ne
resserre le relâché, et l'oracle n'a plus le temps de conclure. Ensuite parce
que la mesure le confirme brutalement : poser $40$ coupes d'archive
\emph{en plus} des $40$ de dominance fait tomber l'optimalité prouvée de
$84$ à $56$ sur $90$ instances.

\medskip\noindent
La règle correcte est donc : \textbf{un seul vivier, un seul plafond, les
deux sources en concurrence}, classées par valeur décroissante du substitut
linéaire $w^{\top}x$ où $w=Qc-Pd$. Ce sont les points de plus grande valeur
du substitut qui tirent $U$ vers le haut, quelle que soit leur origine ;
les couper est ce qui resserre le plus.

L'effet est automatique et souhaitable : là où les points dominés abondent,
ils occupent le vivier ; là où ils sont rares -- le régime à $\Ecal$ épais,
justement -- l'archive le remplit à leur place. Aucun réglage n'arbitre :
le classement le fait.

% ============================================================================
\section{L'algorithme}
\label{sec:algo}

\subsection{Vue d'ensemble}

\begin{figure}[htbp]
\centering
\resizebox{0.99\textwidth}{!}{%
\begin{tikzpicture}[
  font=\small, x=1cm, y=1cm,
  bloc/.style={rectangle,draw,rounded corners=2pt,align=center,
               inner sep=4pt,minimum height=8mm,text width=56mm},
  blocA/.style={bloc,fill=cfonddom,draw=cdom},
  blocB/.style={bloc,fill=cfond,draw=ceff},
  dec/.style={diamond,draw,aspect=2.6,align=center,inner sep=1pt,
              font=\footnotesize,fill=white,text width=25mm},
  fin/.style={rectangle,draw,rounded corners=8pt,align=center,fill=cfond,
              inner sep=4pt,text width=38mm},
  fl/.style={-{Stealth[length=2mm]}},
  et/.style={font=\scriptsize,inner sep=1.5pt,fill=white}
]

\node[bloc,fill=black!4]        (in)      at (0,  0) {instance, budget $T$};
\node[blocA]                    (rech)    at (0, -2.1)
  {\textbf{Phase 1 -- recherche}\\[1pt]
   \footnotesize VNS dans l'espace des critères ; chaque sous-problème
   résolu \emph{exactement},\\
   \footnotesize chaque incumbent \emph{certifié efficace}};
\node[bloc,fill=black!4,text width=96mm] (sortie1) at (0, -4.4)
  {sorties de la phase~1 :\ \ $q=f(\bar x)$ borne inférieure \emph{valide}
   \ \ $\bullet$\ \ archive $\Acal\subseteq\Ecal$
   \ \ $\bullet$\ \ dominés $\mathcal{D}$};
\node[blocB,text width=66mm]    (vivier)  at (0, -6.3)
  {\textbf{vivier unique} (\S\ref{sec:vivier})\\[1pt]
   \footnotesize $\mathcal{D}\cup\Acal$ trié par $w^{\top}x$ décroissant,
   sous un plafond global};
\node[dec]                      (src)     at (0, -8.4) {candidat\\ issu de $\Acal$ ?};
\node[dec]                      (clot)    at (0,-10.6) {(C$_a$)\\ infaisable ?};
\node[blocA,text width=46mm]    (maj)     at (-7.4,-10.6)
  {$q\leftarrow f(x)$, archive enrichie\\[-1pt]
   \footnotesize $x$ est efficace et bat $q$};
\node[blocB,text width=62mm]    (pose)    at (0,-12.9)
  {poser la coupe~\eqref{eq:disj} en $a$\\[-1pt]
   \footnotesize dominance si $a\in\mathcal{D}$,
   \textbf{efficacité} si $a\in\Acal$};
\node[blocB,text width=66mm]    (oracle)  at (0,-15.0)
  {\textbf{Phase 2 -- oracle}\\[1pt]
   \footnotesize maximiser $w^{\top}x+w_0$ sur $\Rcal_A$, à budget borné};
\node[dec]                      (vide)    at (0,-17.0) {$\Rcal_A=\emptyset$ ?};
\node[fin]                      (prop)    at (6.6,-17.0)
  {\textbf{$\qs=q$}\\[-1pt]\footnotesize Proposition~\ref{prop:vide}};
\node[dec]                      (u0)      at (0,-19.0) {$U\le 0$ ?};
\node[fin]                      (prv)     at (6.6,-19.0)
  {\textbf{$\qs=q$}\\[-1pt]\footnotesize borne~\eqref{eq:borne}};
\node[fin,text width=54mm]      (out)     at (0,-21.2)
  {$q\ \le\ \qs\ \le\ q+\dfrac{U}{Q\,D^{+}}$\\[-1pt]
   \footnotesize Théorème~\ref{th:archbound}};

\draw[fl] (in)      -- (rech);
\draw[fl] (rech)    -- (sortie1);
\draw[fl] (sortie1) -- (vivier);
\draw[fl] (vivier)  -- (src);
\draw[fl] (src)     -- node[et,right]{oui} (clot);
\draw[fl] (clot)    -- node[et,right]{oui} (pose);
\draw[fl] (clot)    -- node[et,above]{non} (maj);
\draw[fl] (src.east) -- ++(1.9,0) |- node[et,pos=0.78,above]{non} (pose.east);
\draw[fl] (maj.west) -- ++(-1.3,0)
    |- node[et,pos=0.80,above]{recommencer le tour} (vivier.west);
\draw[fl] (pose)    -- (oracle);
\draw[fl] (oracle)  -- (vide);
\draw[fl] (vide)    -- node[et,above]{oui} (prop);
\draw[fl] (vide)    -- node[et,right]{non} (u0);
\draw[fl] (u0)      -- node[et,above]{oui} (prv);
\draw[fl] (u0)      -- node[et,right]{non} (out);
\draw[fl,dashed,cgris] (out.east) -- ++(3.0,0)
    |- node[et,pos=0.80,above,text=black]{tour suivant} (vivier.east);
\end{tikzpicture}}
\caption{Fonctionnement de l'algorithme. En \textcolor{cdom}{orange} ce qui
existait, en \textcolor{ceff}{bleu} ce que la coupe d'efficacité ajoute. La
flèche décisive est celle qui va de l'archive (sortie de la phase~1) au
vivier de coupes de la phase~2 : c'est l'inversion du flux d'information qui
fait l'hybridation.}
\label{fig:flow}
\end{figure}

\subsection{Pseudo-code}

Trois procédures suffisent. La première établit la clôture, la deuxième
construit le vivier, la troisième est la boucle de certification.

\begin{algorithm}[htbp]
\caption{\textsc{Clôture}$(a,q)$ -- un seul programme de faisabilité}
\label{alg:clot}
\KwIn{$a\in\Ecal$ certifié, incumbent $q=P/Q$}
\KwOut{\textsc{Établie}, ou un point $x$ \emph{efficace} avec $f(x)>q$}
Construire les lignes $Ax\le b$, $x\ge 0$ entier\;
\For{$k\leftarrow 1$ \KwTo $p$}{
  ajouter $e^a_k(x)=0$
  \Comment*[r]{$\iff \Zk(x)=\Zk(a)$, exact car $e^a_k\in\mathbb{Z}$}
}
ajouter $Q\,N(x)-P\,D(x)\ \ge\ 1$
\Comment*[r]{$\iff f(x)>q$, exact car entier}
$r\leftarrow$ résoudre ce programme de \emph{faisabilité}\;
\lIf{$r$ infaisable}{\Return \textsc{Établie}}
\Return $r.x$ \Comment*[r]{efficace \emph{et} meilleur que $q$}
\end{algorithm}

\begin{algorithm}[htbp]
\caption{\textsc{Vivier}$(\mathcal{D},\Acal,w)$ -- une seule réserve, deux sources}
\label{alg:vivier}
\KwIn{points dominés $\mathcal{D}$, archive $\Acal$, substitut $w=Qc-Pd$}
\KwOut{liste de couples (point, origine) triée}
$V\leftarrow\emptyset$\;
\ForEach{$x\in\mathcal{D}$ \emph{distincts}}{$V\leftarrow V\cup\{(x,\textsf{dom})\}$}
\ForEach{$a\in\Acal$ \emph{distincts}}{$V\leftarrow V\cup\{(a,\textsf{arch})\}$}
trier $V$ par $w^{\top}x$ \emph{décroissant}
\Comment*[r]{ce sont eux qui tirent $U$ vers le haut}
\Return $V$
\end{algorithm}

\begin{algorithm}[htbp]
\caption{\textsc{Certifier}$(q,\bar x,\mathcal{D},\Acal,\text{budget})$ --
la boucle de certification}
\label{alg:cert}
\KwIn{incumbent $q$ atteint en $\bar x\in\Ecal$, $\mathcal{D}$, $\Acal$,
      budget $B$, taille de lot $\beta$, nombre de tours $\rho$}
\KwOut{$q_{\text{lb}}\le\qs\le q_{\text{ub}}$, et le drapeau \textsf{prouvé}}
$\Rcal\leftarrow\Scal$ ; $q_{\text{ub}}\leftarrow+\infty$ ;
$n_{\text{arch}}\leftarrow 0$\;
$V\leftarrow$ \textsc{Vivier}$(\mathcal{D},\Acal,w(q))$\;
\For{$t\leftarrow 1$ \KwTo $\rho$}{
  \lIf{budget épuisé}{\Break}
  $q_{\text{réf}}\leftarrow q$ ; $(w,w_0,P,Q)\leftarrow$ substitut$(q_{\text{réf}})$\;
  \tcp{— un lot de coupes de plus, les deux sources mêlées —}
  \ForEach{$(a,\text{origine})$ \emph{des} $\beta$ \emph{premiers de} $V$}{
    \eIf{origine $=\textsf{arch}$}{
      $r\leftarrow$ \textsc{Clôture}$(a,q_{\text{réf}})$\;
      \lIf{$r\ne\textsc{Établie}$}{$q\leftarrow f(r)$ ; $\bar x\leftarrow r$ ;
           $\Acal\leftarrow\Acal\cup\{r\}$ ; \textbf{recommencer le tour}}
      $\Rcal\leftarrow\Rcal\cap\{x:\exists k,\,e^a_k(x)\ge 1\}$ ;
      $n_{\text{arch}}{+}{+}$
      \Comment*[r]{coupe d'\textbf{efficacité}, Th.~\ref{th:effcut}}
    }{
      $\Rcal\leftarrow\Rcal\cap\{x:\exists k,\,e^a_k(x)\ge 1\}$
      \Comment*[r]{coupe de dominance}
    }
    retirer $a$ de $V$\;
  }
  $r\leftarrow$ maximiser $w^{\top}x+w_0$ sur $\Rcal$, à budget borné\;
  \ForEach{point efficace $y$ rendu par l'oracle}{
    $\Acal\leftarrow\Acal\cup\{y\}$ ;
    \lIf{$f(y)>q$}{$q\leftarrow f(y)$ ; $\bar x\leftarrow y$}
  }
  \If{$\Rcal=\emptyset$ \textbf{et} $n_{\text{arch}}>0$}{
    \Return $(q,q,\textsf{prouvé})$
    \Comment*[r]{Proposition~\ref{prop:vide}}
  }
  \lIf{$r$ résolu \textbf{et} $r.\text{valeur}\le 0$}{\Return $(q,q,\textsf{prouvé})$}
  \If{$r.U$ fini}{
    $D^{+}\leftarrow\min\{D(x):x\in\Rcal,\ QN(x)-PD(x)\ge 0\}$\;
    $q_{\text{ub}}\leftarrow\min\bigl(q_{\text{ub}},\
      q_{\text{réf}}+U/(Q\,D^{+})\bigr)$
    \Comment*[r]{Th.~\ref{th:archbound}}
  }
}
\Return $(q,q_{\text{ub}},\ q_{\text{ub}}\le q)$
\end{algorithm}

\begin{remark}[validité malgré un incumbent mobile]
Chaque tour produit une borne valide sur $\qs$ pour le $q$ \emph{de ce tour}.
Comme $\qs$ ne dépend pas de $q$, le \emph{minimum} des bornes obtenues reste
une borne valide : un incumbent amélioré en cours de route ne peut pas
invalider une borne posée plus tôt. C'est pourquoi la ligne du
théorème~\ref{th:archbound} se lit sur $q_{\text{réf}}$ et non sur $q$
courant.
\end{remark}

\subsection{Code}

Les trois pièces réellement nouvelles, en \textsc{Python}. Le solveur entier
est appelé à travers \texttt{scipy.optimize.milp} (HiGHS), et
\texttt{solve\_ilp} en est l'enveloppe qui compte les appels.

\begin{lstlisting}[caption={La condition de clôture : UN seul programme de
faisabilité (Alg.~\ref{alg:clot}).},label={lst:clot}]
def same_criteria_improves(inst, a, q):
    """Existe-t-il x de MEME vecteur criteres que `a` avec f(x) > q ?

    Renvoie None si INFAISABLE  -> cloture etablie, on peut couper en `a`;
    renvoie x    si FAISABLE    -> x est efficace ET meilleur que q.
    """
    P, Q = q.numerator, q.denominator
    rows = feasibility_rows(inst)                  # A x <= b, x >= 0
    for k in range(inst.p):
        coef, cst = e_row(inst, a, k)
        rows.append((coef, -cst, -cst))            # e_k^a(x) = 0
    w  = (Q * inst.f.num - P * inst.f.den).astype(float)
    w0 = float(Q * inst.f.a - P * inst.f.b)
    rows.append((w, 1.0 - w0, INF))                # Q N(x) - P D(x) >= 1
    res = solve_ilp(np.zeros(inst.n), rows, inst.var_upper_bounds(),
                    maximize=True)
    return None if not res.ok else res.x
\end{lstlisting}

\begin{lstlisting}[caption={Le vivier unique : deux sources, un plafond
(Alg.~\ref{alg:vivier}).},label={lst:vivier}]
def build_cut_pool(dominated, archive, w):
    """Les deux coupes ont la meme forme, donc le meme cout (p binaires).
    On les met en concurrence dans UNE seule reserve, triee par valeur
    decroissante du substitut, sous le plafond global deja regle.
    """
    wv = np.asarray(w, dtype=float)
    seen, pool = set(), []
    for pts, kind in ((dominated, "dom"), (archive or [], "arch")):
        for x in pts:
            key = (kind, tuple(int(v) for v in x))
            if key in seen:
                continue
            seen.add(key)
            pool.append((np.asarray(x, dtype=int), kind))
    pool.sort(key=lambda t: -float(wv @ t[0]))
    return pool
\end{lstlisting}

\begin{lstlisting}[caption={Le c\oe ur de la boucle de certification : pose
des coupes, puis les deux issues de preuve (Alg.~\ref{alg:cert}).},
label={lst:cert}]
# --- pose d'un lot de coupes, les deux sources melees -------------------
for x, kind in pending:
    if posees >= cut_batch or time.time() > t0 + budget:
        break
    if kind == "arch":
        better = same_criteria_improves(inst, x, q_ref)   # cloture, Th. 2
        if better is not None:                            # issue FAISABLE
            improved_by_closure = better                  # -> le LB monte
            break
        model.add_efficiency_cut(x)                       # Th. 1
        info["archive_cuts"] += 1
    else:
        model.add_dominance_cut(x)                        # brique 3
    posees += 1

r = max_linear_over_E(inst, w, w0, time_limit=slice_, model=model)

# --- RELACHE VIDE : la preuve la plus forte -----------------------------
# Th. 1 : E est inclus dans R_A union {x : Z(x) = Z(a), a dans A}.
# Si R_A est vide, tout point efficace a le vecteur criteres d'un point de
# l'archive, et la cloture -- etablie point par point AVANT chaque coupe --
# dit qu'aucun d'eux ne bat q. Donc q* <= q, et comme q <= q*, q* = q.
# La coupe de DOMINANCE ne peut jamais produire cette preuve : elle
# preserve E, donc R ne se vide pas.
if r.status == "empty" and info["archive_cuts"] > 0:
    return CertResult(float(q), True, q, x_cur, info)

# --- optimalite prouvee par la borne ------------------------------------
if r.status == "optimal" and r.value <= 0 and not improved:
    return CertResult(float(q), True, q, x_cur, info)
\end{lstlisting}

% ============================================================================
\section{Illustration numérique 1 : deux variables, pas à pas}
\label{sec:illus2}

\subsection{L'instance}

\[
\mathcal{I}_2:\quad
\begin{aligned}
  &\max\ Z_1(x)=\frac{3x_2+2}{x_1+x_2+3},
   \qquad Z_2(x)=\frac{x_1+3x_2}{2x_2+3},\\[3pt]
  &\max\ f(x)=\frac{3x_1+x_2}{2x_2+1}
   \quad\text{sur }\Ecal,\\[3pt]
  &\text{s.c. } 3x_1+x_2\le 14,\quad 2x_1+3x_2\le 15,
   \quad x\in\mathbb{Z}^2_+ .
\end{aligned}
\]

\noindent
Les hypothèses (A1)--(A2) sont vérifiées par inspection. L'énumération
exhaustive donne $|\Scal|=22$, $|\Ecal|=5$ et $\qs=14/5=2{,}8$, atteint en
$x^{\star}=(4,2)$. Ces valeurs ne servent qu'à \emph{vérifier} ce que la
méthode produit ; elle ne les utilise jamais.

\begin{remark}[reproductibilité]
Tous les chiffres des sections~\ref{sec:illus2} et~\ref{sec:illus3} --
tableaux, trajectoires de borne, coordonnées des figures -- sont produits par
le script \texttt{illustration.py}, qui réexécute la méthode et revérifie
chaque encadrement. Aucun n'est recopié à la main.
\end{remark}

\begin{figure}[htbp]
\centering
\begin{tikzpicture}[scale=1.32,font=\small]
% polyedre
\fill[cfond,opacity=0.75] (0,0) -- (4.667,0) -- (3.857,2.429) -- (0,5) -- cycle;
\draw[cgris] (0,0) -- (4.667,0) -- (3.857,2.429) -- (0,5) -- cycle;
% axes
\draw[-{Stealth[length=2mm]}] (-0.25,0) -- (5.2,0) node[right]{$x_1$};
\draw[-{Stealth[length=2mm]}] (0,-0.25) -- (0,5.6) node[above]{$x_2$};
\foreach \i in {0,1,2,3,4,5} {
  \draw (\i,-0.06) -- (\i,0.06) node[below=3pt,font=\scriptsize]{\i};
  \draw (-0.06,\i) -- (0.06,\i) node[left=3pt,font=\scriptsize]{\i};
}
% points dominés
\foreach \p in {(0,0),(0,1),(0,2),(0,3),(0,4),(1,0),(1,1),(1,2),(1,3),
                (2,0),(2,1),(2,2),(3,0),(3,1),(3,2),(4,1)}
  \fill[cgris] \p circle (1.5pt);
% argmax f sur S : (4,0), dominé
\fill[cdom] (4,0) circle (2.4pt);
\draw[cdom,thick] (4,0) circle (5pt);
\node[cdom,font=\scriptsize,fill=white,inner sep=1.5pt,left=7pt]
  at (4,0.34) {$\max_{\Scal} f=12$};
% points efficaces
\foreach \p in {(0,5),(1,4),(2,3),(3,3),(4,2)}
  {\fill[ceff] \p circle (2.6pt);}
\draw[ceff,thick,dashed] (0,5) -- (1,4) -- (2,3) -- (3,3) -- (4,2);
\node[ceff,font=\scriptsize,above right=0pt and 1pt] at (4,2) {$x^\star=(4,2)$};
\node[ceff,font=\scriptsize,above right=0pt and 1pt] at (0,5) {$\Ecal$};
% contraintes
\node[cgris,font=\scriptsize,rotate=-72] at (4.35,1.5) {$3x_1+x_2\le 14$};
\node[cgris,font=\scriptsize,rotate=-33] at (1.9,4.1) {$2x_1+3x_2\le 15$};
\end{tikzpicture}
\caption{$\mathcal{I}_2$ dans l'espace des décisions. Les $22$ points entiers
de $\Scal$ ; en \textcolor{ceff}{bleu} les $5$ points efficaces, en
\textcolor{cdom}{orange} le point $(4,0)$ où $f$ atteint son maximum sur
$\Scal$, égal à $12$. Ce point est \emph{dominé} : borner $\qs$ par
$\max_\Scal f$ donne $12$ au lieu de $2{,}8$, soit un écart de $77\ \%$.
Tout le travail consiste à retrancher ce genre de région.}
\label{fig:i2x}
\end{figure}

\subsection{Étape 1 -- la recherche}

La phase~1 explore l'espace des critères et certifie chaque incumbent par la
brique~1. Sur cette instance elle rend, dans cet ordre :
\[
  (0,5),\quad (2,3),\quad (1,4),\quad (3,3),\quad (4,2),
\]
c'est-à-dire l'archive $\Acal$ des $5$ points efficaces, et l'incumbent
$q=f(4,2)=14/5$. \emph{La réparation atteint un point efficace du premier
coup : la recherche ne collecte aucun point dominé.} Le vivier de la coupe de
dominance est donc \textbf{vide} -- c'est exactement le verrou de
la section~\ref{sec:verrou}, sur un cas concret.

\subsection{Étape 2 -- ce que chaque coupe retranche}

\begin{figure}[htbp]
\centering
\begin{tikzpicture}[scale=3.05,font=\small]
% zone efficacité (a = (4,2), Z = (0.8889, 1.4286))
\fill[cfond] (0.20,-0.02) rectangle (0.8889,1.4286);
% zone dominance (xbar = (4,0), Z = (0.2857, 1.3333))
\fill[cfonddom,opacity=0.95] (0.20,-0.02) rectangle (0.2857,1.3333);
\draw[cdom,thick] (0.2857,-0.02) -- (0.2857,1.3333) -- (0.20,1.3333);
\draw[ceff,thick] (0.8889,-0.02) -- (0.8889,1.4286) -- (0.20,1.4286);
% axes
\draw[-{Stealth[length=2mm]}] (0.20,0) -- (2.28,0) node[right]{$Z_1$};
\draw[-{Stealth[length=2mm]}] (0.22,-0.02) -- (0.22,1.58) node[above]{$Z_2$};
\foreach \v in {0.5,1.0,1.5,2.0}
  \draw (\v,-0.015) -- (\v,0.015) node[below=2pt,font=\scriptsize]{\v};
\foreach \v in {0.0,0.5,1.0,1.5}
  \draw (0.205,\v) -- (0.235,\v) node[left=2pt,font=\scriptsize]{\v};
% points dominés
\foreach \p in {(0.6667,0.0),(1.25,0.6),(1.6,0.8571),(1.8333,1.0),(2.0,1.0909),
                (0.5,0.3333),(1.0,0.8),(1.3333,1.0),(1.5714,1.1111),
                (0.4,0.6667),(0.8333,1.0),(1.1429,1.1429),
                (0.3333,1.0),(0.7143,1.2),(1.0,1.2857),(0.625,1.4)}
  \fill[cgris] \p circle (0.55pt);
\fill[cdom] (0.2857,1.3333) circle (1.1pt);
\node[cdom,font=\scriptsize,below right=0pt and 1pt] at (0.2857,1.3333)
  {$\bar x=(4,0)$};
\node[ceff,font=\scriptsize,align=center] at (0.60,0.72)
  {\textbf{9 points}\\ retranchés};
\node[cdom,font=\scriptsize,align=center,fill=white,inner sep=1pt]
  at (0.35,0.30) {};
\draw[cdom,thin] (0.2650,0.36) -- (0.44,0.20);
\node[cdom,font=\scriptsize,right=0pt] at (0.44,0.185)
  {\textbf{1 point} retranché};
% front
\foreach \p in {(2.125,1.1538),(1.75,1.1818),(1.375,1.2222),(1.2222,1.3333),(0.8889,1.4286)}
  \fill[ceff] \p circle (1.1pt);
\draw[ceff,dashed] (2.125,1.1538) -- (1.75,1.1818) -- (1.375,1.2222)
     -- (1.2222,1.3333) -- (0.8889,1.4286);
\node[ceff,font=\scriptsize,above right=0pt and 1pt] at (0.8889,1.4286) {$a=(4,2)$};
\node[ceff,font=\scriptsize,right=2pt] at (2.125,1.1538) {$(0,5)$};
\end{tikzpicture}
\caption{Espace des critères de $\mathcal{I}_2$ : les $22$ images, le front
en \textcolor{ceff}{bleu}. La coupe de \textcolor{cdom}{dominance} en
$\bar x=(4,0)$ retranche le rectangle orange -- et ce rectangle ne contient
qu'\textbf{un seul} point, $\bar x$ lui-même : $|\Rcal|$ passe de $22$ à
$21$. La coupe d'\textcolor{ceff}{efficacité} en $a=(4,2)$ retranche le
rectangle bleu, qui en contient \textbf{neuf} : $|\Rcal_A|$ passe de $22$ à
$13$. Même forme algébrique, même coût, neuf fois plus de matière.}
\label{fig:i2z}
\end{figure}

\subsection{Étape 3 -- la clôture}

Avant chaque coupe d'efficacité on résout (C$_a$) de~\eqref{eq:clot}. Ici,
pour les $5$ points de l'archive et $q=14/5$ :
\begin{center}
\begin{tabular}{lccccc}
\toprule
$a$ & $(4,2)$ & $(3,3)$ & $(2,3)$ & $(1,4)$ & $(0,5)$\\
\midrule
(C$_a$) & infaisable & infaisable & infaisable & infaisable & infaisable\\
verdict & \multicolumn{5}{c}{clôture établie -- les $5$ coupes sont licites}\\
\bottomrule
\end{tabular}
\end{center}
\noindent
Coût total : $5$ appels au solveur entier, un par point. À comparer aux $5$
appels de Dinkelbach \emph{complets} qu'exigerait une lecture littérale
de~\eqref{eq:cloture}.

\subsection{Étape 4 -- le relâché se vide}

Le vivier est trié par $w^{\top}x$ décroissant, avec $w=Qc-Pd=(15,-23)$ pour
$q=14/5$. L'ordre est donc $(4,2),(3,3),(2,3),(1,4),(0,5)$. La
figure~\ref{fig:i2panels} montre $\Rcal_A$ après chaque coupe.

\newcommand{\panel}[2]{%
\begin{tikzpicture}[scale=0.46,font=\tiny]
  \draw[cgris!50,very thin] (-0.45,-0.45) rectangle (4.55,5.55);
  \foreach \p in {(0,0),(0,1),(0,2),(0,3),(0,4),(0,5),(1,0),(1,1),(1,2),(1,3),
                  (1,4),(2,0),(2,1),(2,2),(2,3),(3,0),(3,1),(3,2),(3,3),
                  (4,0),(4,1),(4,2)}
    {\draw[cgris] \p circle (2.0pt);}
  \foreach \p in {#2} \fill[ceff] \p circle (3.4pt);
  \node[below=2pt,font=\scriptsize,align=center] at (2.05,-0.75) {#1};
\end{tikzpicture}}

\begin{figure}[htbp]
\centering
\begin{tabular}{@{}c@{\hspace{9mm}}c@{\hspace{9mm}}c@{}}
\panel{$0$ coupe\\$|\Rcal_A|=22$}{(0,0),(0,1),(0,2),(0,3),(0,4),(0,5),(1,0),(1,1),(1,2),(1,3),(1,4),(2,0),(2,1),(2,2),(2,3),(3,0),(3,1),(3,2),(3,3),(4,0),(4,1),(4,2)}
&
\panel{coupe en $(4,2)$\\$|\Rcal_A|=13$}{(0,1),(0,2),(0,3),(0,4),(0,5),(1,1),(1,2),(1,3),(1,4),(2,2),(2,3),(3,2),(3,3)}
&
\panel{coupe en $(3,3)$\\$|\Rcal_A|=9$}{(0,1),(0,2),(0,3),(0,4),(0,5),(1,2),(1,3),(1,4),(2,3)}
\\[7mm]
\panel{coupe en $(2,3)$\\$|\Rcal_A|=6$}{(0,2),(0,3),(0,4),(0,5),(1,3),(1,4)}
&
\panel{coupe en $(1,4)$\\$|\Rcal_A|=3$}{(0,3),(0,4),(0,5)}
&
\panel{coupe en $(0,5)$\\$\Rcal_A=\emptyset$ \textbf{: preuve}}{}
\end{tabular}
\caption{Les six états successifs du relâché $\Rcal_A$ de $\mathcal{I}_2$
(espace des décisions, $x_1$ en abscisse). Disques pleins : $\Rcal_A$ ;
cercles vides : retranchés. Après les cinq coupes d'efficacité le relâché est
\textbf{vide}, et la proposition~\ref{prop:vide} conclut $\qs=q=14/5$
\emph{sans qu'aucune borne n'ait été évaluée}. Aucune suite de coupes de
dominance ne peut atteindre cet état : elles préservent $\Ecal$, donc
$|\Rcal|\ge 5$ toujours.}
\label{fig:i2panels}
\end{figure}

\begin{figure}[htbp]
\centering
\begin{tikzpicture}
\begin{axis}[width=11.6cm,height=5.4cm,
  xlabel={nombre de coupes posées}, ylabel={$|\Rcal_A|$},
  xmin=-0.3,xmax=5.4, ymin=-2, ymax=40,
  xtick={0,1,2,3,4,5}, ytick={0,10,20,30},
  legend style={at={(0.98,0.96)},anchor=north east,font=\small,draw=cgris},
  grid=major, grid style={cgris!35}, tick label style={font=\small},
  label style={font=\small},
  nodes near coords, nodes near coords style={font=\scriptsize}]
\addplot[color=ceff,mark=*,thick,mark size=2.4pt]
  coordinates {(0,22) (1,13) (2,9) (3,6) (4,3) (5,0)};
\addlegendentry{$\mathcal{I}_2$ : coupes d'efficacité}
\addplot[color=black!55,mark=triangle*,thick,dashed,mark size=2.6pt]
  coordinates {(0,36) (1,23) (2,16) (3,5) (4,1)};
\addlegendentry{$\mathcal{I}_3$ : coupes d'efficacité}
\addplot[color=cdom,mark=square*,thick,dotted,mark size=2.2pt]
  coordinates {(0,22) (1,21) (2,20) (3,17) (4,17) (5,15)};
\addlegendentry{$\mathcal{I}_2$ : coupes de dominance}
\end{axis}
\end{tikzpicture}
\caption{Le relâché, coupe après coupe. Sur $\mathcal{I}_2$ la coupe
d'efficacité le fait fondre de $22$ points à \textbf{zéro} en cinq coupes --
et un relâché vide \emph{est} la preuve (proposition~\ref{prop:vide}). La
coupe de dominance, sur la même instance et pour le même nombre de coupes,
descend de $22$ à $15$ : elle ne peut pas faire mieux, puisqu'elle préserve
$\Ecal$ et se heurte donc au plancher $|\Rcal|\ge|\Ecal|=5$. Sur
$\mathcal{I}_3$ la descente s'arrête à $1$ : l'archive y est incomplète, il
reste le point efficace non trouvé -- et le théorème~\ref{th:archbound}
conclut malgré tout, dès la deuxième coupe.}
\label{fig:trajrelache}
\end{figure}

\subsection{Étape 5 -- la borne, coupe par coupe}

Le tableau~\ref{tab:i2traj} donne, à chaque étape, les quantités exactes du
théorème~\ref{th:archbound}. \emph{Une seule} coupe d'efficacité suffit à
faire tomber $U$ à $0$, donc à prouver l'optimalité.

\begin{table}[htbp]
\centering
\caption{$\mathcal{I}_2$ : trajectoire des deux coupes, $q=14/5$,
$U=\max_{\Rcal}\{QN-PD\}$, borne $=q+U/(QD^{+})$. Les valeurs sont exactes
(rationnelles).}
\label{tab:i2traj}
\setlength{\tabcolsep}{6pt}
\begin{tabular}{clccc@{\qquad}lccc}
\toprule
& \multicolumn{4}{c}{\textbf{coupe d'efficacité} (source : archive)}
& \multicolumn{4}{c}{coupe de dominance (source : réparations)}\\
\cmidrule(lr){2-5}\cmidrule(lr){6-9}
\# & coupe en & $|\Rcal|$ & $U$ & borne
   & coupe en & $|\Rcal|$ & $U$ & borne\\
\midrule
0 & --        & 22 & 46 & $12$
  & --        & 22 & 46 & $12$\\
1 & $(4,2)$   & 13 & $\mathbf{0}$ & $\mathbf{14/5}$
  & $(4,0)$   & 21 & 31 & $9$\\
2 & $(3,3)$   &  9 & 0 & $14/5$
  & $(3,0)$   & 20 & 23 & $37/5$\\
3 & $(2,3)$   &  6 & 0 & $14/5$
  & $(4,1)$   & 17 &  8 & $10/3$\\
4 & $(1,4)$   &  3 & 0 & $14/5$
  & $(2,0)$   & 17 &  8 & $10/3$\\
5 & $(0,5)$   & $\mathbf{0}$ & \multicolumn{2}{c}{\textbf{vide} $\Rightarrow \qs=q$}
  & $(3,1)$   & 15 & $0$ & $14/5$\\
\midrule
\multicolumn{5}{l}{\textbf{disponibles : 5 (l'archive)}}
& \multicolumn{4}{l}{disponibles : \textbf{0} (aucun dominé collecté)}\\
\bottomrule
\end{tabular}
\end{table}

\begin{figure}[htbp]
\centering
\begin{tikzpicture}
\begin{axis}[width=12.4cm,height=6.2cm,
  xlabel={nombre de coupes posées}, ylabel={borne supérieure sur $\qs$},
  xmin=-0.3,xmax=9.3, ymin=2.2, ymax=12.8,
  xtick={0,1,2,3,4,5,6,7,8,9},
  legend style={at={(0.98,0.96)},anchor=north east,font=\small,draw=cgris},
  grid=major, grid style={cgris!35}, tick label style={font=\small},
  label style={font=\small}]
\addplot[color=ceff,mark=*,thick,mark size=2pt]
  coordinates {(0,12) (1,2.8) (2,2.8) (3,2.8) (4,2.8) (5,2.8)};
\addlegendentry{coupe d'efficacité (archive)}
\addplot[color=cdom,mark=square*,thick,mark size=1.9pt,dashed]
  coordinates {(0,12) (1,9) (2,7.4) (3,3.3333) (4,3.3333) (5,2.8)
               (6,2.8) (7,2.8) (8,2.8) (9,2.8)};
\addlegendentry{coupe de dominance (17 points, hypothétiques)}
\addplot[color=black,thick,dotted,domain=-0.3:9.3] {2.8};
\addlegendentry{$\qs=14/5$}
\end{axis}
\end{tikzpicture}
\caption{$\mathcal{I}_2$ : la borne en fonction du nombre de coupes.
L'\textcolor{ceff}{efficacité} atteint l'optimum en \textbf{une} coupe ; la
\textcolor{cdom}{dominance} en demande \textbf{cinq} -- et elle les prend
dans un vivier de $17$ points que la recherche, ici, n'a \emph{pas produit}.
La courbe orange est donc une borne de ce que la coupe classique pourrait
faire au mieux, pas de ce qu'elle fait.}
\label{fig:i2traj}
\end{figure}

\subsection{Ce que la méthode rend réellement}

Exécution complète sur $\mathcal{I}_2$, budget $3+2$~s :
\begin{center}
\begin{tabular}{lcccc}
\toprule
& $q_{\text{lb}}$ & $q_{\text{ub}}$ & prouvé & coupes d'efficacité posées\\
\midrule
sans coupes d'archive & $14/5$ & $2{,}80$ & oui & $0$\\
avec coupes d'archive & $14/5$ & $2{,}80$ & oui & $5$\\
\bottomrule
\end{tabular}
\end{center}
\noindent
Sur $22$ points, les deux variantes concluent : l'instance est trop petite
pour départager. L'illustration sert à exposer le \emph{mécanisme} ; l'effet
se mesure à la section~\ref{sec:res}, sur $90$ instances.

% ============================================================================
\section{Illustration numérique 2 : trois variables}
\label{sec:illus3}

\subsection{L'instance}

\[
\mathcal{I}_3:\quad
\begin{aligned}
  &\max\ Z_1(x)=\frac{3x_3}{2x_2+x_3+2},
   \qquad Z_2(x)=\frac{3x_1+3x_2+x_3+1}{x_1+2x_2+x_3+3},\\[3pt]
  &\max\ f(x)=\frac{2x_1+2x_2+x_3+1}{x_1+x_3+1}\quad\text{sur }\Ecal,\\[3pt]
  &\text{s.c. } 3x_1+x_2+x_3\le 8,\quad x_1+2x_2+2x_3\le 9,
   \quad x\in\mathbb{Z}^3_+ .
\end{aligned}
\]

\noindent
$|\Scal|=36$, $|\Ecal|=5$, $\qs=3$ atteint en $x^{\star}=(2,2,0)$, tandis que
$\max_{\Scal}f=9$ en $(0,4,0)$ -- dominé. L'écart de la borne naïve est de
$67\ \%$. L'archive rendue par la recherche compte $4$ points :
$(1,0,4)$, $(2,0,2)$, $(2,2,0)$, $(2,1,1)$ ; le cinquième point efficace,
$(2,0,1)$, n'est \emph{pas} trouvé. Les points dominés collectés : $0$,
là encore.

\begin{figure}[htbp]
\centering
\begin{tikzpicture}[scale=3.6,font=\small]
\fill[cfond] (-0.03,0.28) rectangle (1.0,1.4444);
\draw[ceff,thick] (1.0,0.28) -- (1.0,1.4444) -- (-0.03,1.4444);
\draw[-{Stealth[length=2mm]}] (-0.03,0.3) -- (2.22,0.3) node[right]{$Z_1$};
\draw[-{Stealth[length=2mm]}] (0,0.27) -- (0,1.60) node[above]{$Z_2$};
\foreach \v in {0.5,1.0,1.5,2.0}
  \draw (\v,0.285) -- (\v,0.315) node[below=2pt,font=\scriptsize]{\v};
\foreach \v in {0.5,1.0,1.5}
  \draw (-0.012,\v) -- (0.012,\v) node[left=2pt,font=\scriptsize]{\v};
% dominés
\foreach \p in {(0,0.3333),(1,0.5),(1.5,0.6),(1.8,0.6667),(2,0.7143),
 (0,0.8),(0.6,0.8333),(1,0.8571),(1.2857,0.875),(0,1),(0.4286,1),(0.75,1),
 (0,1.1111),(0.3333,1.1),(0,1.1818),(0,1),(1,1),(1.5,1),(1.8,1),
 (0,1.1667),(0.6,1.1429),(1,1.125),(1.2857,1.1111),(0,1.25),(0.4286,1.2222),
 (0.75,1.2),(0,1.3),(0.3333,1.2727),(0,1.3333),(0,1.4),(0,1.4286)}
  \fill[cgris] \p circle (0.7pt);
% efficaces
\foreach \p in {(2,1),(1.5,1.2857),(0.6,1.375),(0,1.4444)}
  \fill[ceff] \p circle (1.3pt);
\fill[white] (1,1.3333) circle (1.3pt);
\draw[ceff,thick] (1,1.3333) circle (1.3pt);
\draw[ceff,dashed] (2,1) -- (1.5,1.2857) -- (1,1.3333) -- (0.6,1.375) -- (0,1.4444);
\node[ceff,font=\scriptsize,right=2pt] at (2,1) {$(1,0,4)$};
\node[ceff,font=\scriptsize,above right=0pt and 1pt] at (1.5,1.2857) {$(2,0,2)$};
\node[font=\scriptsize,below=4pt,fill=white,inner sep=1pt] at (1,1.3333)
  {$(2,0,1)$ \emph{non trouvé}};
\node[ceff,font=\scriptsize,above right=0pt and 2pt] at (0.6,1.375) {$(2,1,1)$};
\node[ceff,font=\scriptsize,above right=0pt and 2pt] at (0,1.4444)
  {$x^\star=(2,2,0)$};
\node[ceff,font=\scriptsize,align=center] at (0.47,0.80)
  {zone retranchée\\ par la coupe en $(2,0,1)$};
\end{tikzpicture}
\caption{$\mathcal{I}_3$ dans l'espace des critères (qui reste plan puisque
$p=2$, même à $n=3$). Le point efficace $(2,0,1)$, en cercle vide, n'est pas
dans l'archive. La zone bleue est ce que retranche la coupe d'efficacité en
$(2,0,1)$\dots\ qu'on ne peut pas poser. C'est pourquoi $\Rcal_A$ ne se vide
pas ici : il finit \emph{exactement} sur ce point.}
\label{fig:i3z}
\end{figure}

\subsection{Trajectoire}

\begin{table}[htbp]
\centering
\caption{$\mathcal{I}_3$ : trajectoire, $q=3$. La borne atteint $\qs$ après
\textbf{deux} coupes d'efficacité, contre \textbf{huit} coupes de dominance
-- prises, une fois de plus, dans un vivier que la recherche n'a pas produit.}
\label{tab:i3traj}
\setlength{\tabcolsep}{6pt}
\begin{tabular}{clccc@{\qquad}lccc}
\toprule
& \multicolumn{4}{c}{\textbf{coupe d'efficacité} (archive, 4 points)}
& \multicolumn{4}{c}{coupe de dominance (31 points, hypothétiques)}\\
\cmidrule(lr){2-5}\cmidrule(lr){6-9}
\# & coupe en & $|\Rcal|$ & $U$ & borne
   & coupe en & $|\Rcal|$ & $U$ & borne\\
\midrule
0 & --          & 36 & 6 & $9$      & --          & 36 & 6 & $9$\\
1 & $(2,2,0)$   & 23 & 2 & $4$      & $(0,4,0)$   & 29 & 5 & $11/2$\\
2 & $(2,1,1)$   & 16 & $\mathbf{0}$ & $\mathbf{3}$
                                    & $(1,4,0)$   & 26 & 2 & $4$\\
3 & $(2,0,2)$   &  5 & 0 & $3$      & $(0,3,0)$   & 26 & 2 & $4$\\
4 & $(1,0,4)$   &  1 & 0 & $3$      & $(1,3,0)$   & 26 & 2 & $4$\\
5 & \multicolumn{4}{l}{archive épuisée ; $\Rcal_A=\{(2,0,1)\}$}
  & $(0,2,0)$   & 26 & 2 & $4$\\
6 & & & & & $(0,3,1)$ & 25 & 1 & $7/2$\\
7 & & & & & $(1,2,0)$ & 25 & 1 & $7/2$\\
8 & & & & & $(1,3,1)$ & 24 & $\mathbf{0}$ & $\mathbf{3}$\\
\bottomrule
\end{tabular}
\end{table}

\begin{figure}[htbp]
\centering
\begin{tikzpicture}
\begin{axis}[width=12.4cm,height=5.8cm,
  xlabel={nombre de coupes posées}, ylabel={borne supérieure sur $\qs$},
  xmin=-0.3,xmax=8.3, ymin=2.6, ymax=9.6,
  xtick={0,1,2,3,4,5,6,7,8},
  legend style={at={(0.98,0.96)},anchor=north east,font=\small,draw=cgris},
  grid=major, grid style={cgris!35}, tick label style={font=\small},
  label style={font=\small}]
\addplot[color=ceff,mark=*,thick,mark size=2pt]
  coordinates {(0,9) (1,4) (2,3) (3,3) (4,3)};
\addlegendentry{coupe d'efficacité (4 points d'archive)}
\addplot[color=cdom,mark=square*,thick,mark size=1.9pt,dashed]
  coordinates {(0,9) (1,5.5) (2,4) (3,4) (4,4) (5,4) (6,3.5) (7,3.5) (8,3)};
\addlegendentry{coupe de dominance (31 points, hypothétiques)}
\addplot[color=black,thick,dotted,domain=-0.3:8.3] {3};
\addlegendentry{$\qs=3$}
\end{axis}
\end{tikzpicture}
\caption{$\mathcal{I}_3$ : même lecture qu'à la figure~\ref{fig:i2traj}.
L'écart de coût est ici d'un facteur $4$ en nombre de coupes -- donc en
variables binaires ajoutées au programme entier.}
\label{fig:i3traj}
\end{figure}

\noindent
\textbf{Ce que cette seconde illustration ajoute.} Elle montre le cas où
$\Rcal_A$ \emph{ne} se vide \emph{pas} : l'archive est incomplète, il y
subsiste le point efficace non trouvé. La proposition~\ref{prop:vide} ne
s'applique donc pas -- et ce n'est pas grave, car le
théorème~\ref{th:archbound} conclut quand même, dès la deuxième coupe, par
$U=0$. Les deux voies de preuve sont indépendantes, et c'est ce qui rend la
méthode robuste à la qualité de l'archive.

% ============================================================================
\section{Résultats}
\label{sec:res}

\subsection{Protocole}

Deux jeux d'instances, générés sous (A1)--(A2), gelés avant toute mesure :
un \emph{régime cible} de $8$ instances à vérité terrain énumérable, jouées
sur $10$ graines, et un \emph{lot systématique} de $90$ instances. Les deux
variantes reçoivent un budget \emph{identique} ; l'unité de coût rapportée
est le nombre d'appels au solveur entier, indépendante de la machine. Toutes
les exécutions sont soumises au contrôle
$q_{\text{lb}}\le\qs\le q_{\text{ub}}$ ainsi qu'à quatre vérifications
croisées inter-méthodes.

\subsection{Régime cible}

\begin{table}[htbp]
\centering
\caption{Régime cible, $8$ instances $\times$ $10$ graines, budget identique
de $7+5$~s, vérité terrain disponible.}
\label{tab:res1}
\begin{tabular}{lcc}
\toprule
& sans coupes d'archive & avec\\
\midrule
écart garanti médian & $0{,}00\ \%$ & $0{,}00\ \%$\\
optimalité \textbf{prouvée} & $58/80$ exécutions & $\mathbf{78/80}$\\
validité $q_{\text{lb}}\le\qs\le q_{\text{ub}}$ & $80/80$ & $\mathbf{80/80}$\\
\midrule
\multicolumn{3}{l}{\emph{lecture appariée} (même instance, même graine)}\\
preuves gagnées / perdues & \multicolumn{2}{c}{$\mathbf{20}$ / $\mathbf{0}$}\\
delta des appels par paire & \multicolumn{2}{c}{médiane $\mathbf{-12}$,
  total $-1\,579$, $54/80$ en baisse}\\
\bottomrule
\end{tabular}
\end{table}


\begin{figure}[htbp]
\centering
\begin{tikzpicture}
\begin{axis}[width=13cm,height=5.6cm,
  ybar, bar width=7pt, enlarge x limits=0.08,
  symbolic x coords={{n7 p4 c0},{n7 p4 c.25},{n6 p4 c0},{n7 p3 c.25},
                     {n5 p4 c0},{n8 p3 c0},{n7 p4 c.5},{n6 p2 c.25}},
  xtick=data, xticklabel style={font=\scriptsize,rotate=25,anchor=north east},
  ylabel={optimalité prouvée, sur $10$ graines},
  ymin=0, ymax=11.4, ytick={0,2,4,6,8,10},
  nodes near coords, nodes near coords style={font=\tiny},
  legend style={at={(0.5,1.04)},anchor=south,legend columns=2,
                font=\small,draw=cgris},
  grid=major, grid style={cgris!35}, tick label style={font=\small},
  label style={font=\small}]
\addplot[fill=cdom!45,draw=cdom] coordinates
  {({n7 p4 c0},10) ({n7 p4 c.25},0) ({n6 p4 c0},10) ({n7 p3 c.25},10)
   ({n5 p4 c0},0) ({n8 p3 c0},8) ({n7 p4 c.5},10) ({n6 p2 c.25},10)};
\addlegendentry{sans coupes d'archive}
\addplot[fill=ceff!35,draw=ceff] coordinates
  {({n7 p4 c0},10) ({n7 p4 c.25},8) ({n6 p4 c0},10) ({n7 p3 c.25},10)
   ({n5 p4 c0},10) ({n8 p3 c0},10) ({n7 p4 c.5},10) ({n6 p2 c.25},10)};
\addlegendentry{avec}
\end{axis}
\end{tikzpicture}
\caption{Régime cible, instance par instance, $10$ graines chacune. Six
instances sur huit étaient déjà prouvées à toutes les graines : la coupe
d'efficacité ne pouvait rien y ajouter, et n'y retire rien. Le gain est
entièrement porté par les trois autres, et deux d'entre elles changent
d'\emph{état} plutôt que de degré : de \textbf{zéro} preuve sur dix à dix
sur dix, et de zéro à huit. Elles passent de « solution trouvée, optimalité
inconnue » à « optimalité établie » -- ce qui n'est pas une amélioration de
pourcentage mais un changement de ce que la méthode rend.}
\label{fig:cible}
\end{figure}

\noindent
C'est ici que le théorème~\ref{th:hauteur} change le signe du coût. Avant
lui, les preuves supplémentaires s'\emph{achetaient} : la médiane appariée
valait $+10$ appels. Chaque clôture établie par programme linéaire est un
appel entier non dépensé, et le temps ainsi rendu retourne à la
certification, qui ferme davantage de bornes. Les preuves montent
\emph{et} le coût baisse, non l'un malgré l'autre.

Deux instances changent d'\emph{état}, et non seulement de chiffre : l'une
passe de $0/10$ preuves et $17{,}3\ \%$ d'écart garanti à $8/10$ et
$0{,}0\ \%$ ; l'autre de $0/10$ et $6{,}7\ \%$ à $10/10$ et $0{,}0\ \%$.
Elles passent donc de « solution trouvée, optimalité inconnue » à
« optimalité prouvée ». La première est celle qu'il faut $433$~s et $177$
coupes pour prouver par l'oracle exact seul.

\subsection{Lot systématique}

\begin{table}[htbp]
\centering
\caption{Lot figé, $90$ instances, budget identique de $12$~s. Toutes les
vérifications croisées passent.}
\label{tab:res2}
\setlength{\tabcolsep}{6pt}
\begin{tabular}{lccccc}
\toprule
& optimalité & meilleure & \ILP & $t$ & vérifications\\
& \textbf{prouvée} & valeur & méd. & méd. & en échec\\
\midrule
méthode exacte      & $69/90$ & $75/90$ & $92$  & $1{,}8$~s & 0\\
recherche seule     & $84/90$ & $\mathbf{90/90}$ & $124$ & $0{,}9$~s & 0\\
\textbf{+ coupes d'efficacité} & $\mathbf{86/90}$ & $\mathbf{90/90}$
                    & $\mathbf{97}$ & $\mathbf{0{,}7}$~s & 0\\
\bottomrule
\end{tabular}
\end{table}

\noindent
Deux preuves de plus et aucune perdue, pour une qualité de solution
identique : la recherche atteint déjà l'optimum sur les $90$ instances, il
n'y a donc pas de place pour un progrès de ce côté.

\paragraph{Une lecture qui se trompe, et celle qui la corrige.} Les colonnes
$\ILP$ et $t$ du tableau~\ref{tab:res2} sont les médianes de deux
distributions \emph{indépendantes}. Les lire comme un gain par instance est
une faute, et nous l'avons commise avant de la mesurer. La comparaison
correcte est \textbf{appariée} : même instance, même graine, on retranche.

\begin{table}[htbp]
\centering
\caption{Lecture appariée du même lot. Le signe change.}
\label{tab:apparie}
\begin{tabular}{lcc}
\toprule
& lecture par médianes & lecture \textbf{appariée}\\
\midrule
appels au solveur entier & $124\to 97$, soit $-22\ \%$
                         & médiane $\mathbf{+1}$ par paire\\
temps                    & $0{,}9\to 0{,}7$~s
                         & médiane $\mathbf{+0{,}0}$~s\\
instances moins chères   & --- & $40/90$ \ (plus chères : $47$)\\
total des appels         & --- & $-5{,}2\ \%$\\
total du temps           & --- & $-23\ \%$\\
\bottomrule
\end{tabular}
\end{table}

\noindent
Le gain existe donc, mais il est \emph{concentré dans une queue courte} --
la plus forte baisse atteint $205$ appels sur une seule instance -- et non
réparti. L'instance médiane ne gagne rien. Écrire « strictement positif sur
tous les axes », comme le suggérait la première lecture, aurait été faux.


\begin{figure}[htbp]
\centering
\begin{tikzpicture}
\begin{axis}[width=12.6cm,height=6.6cm,
  xlabel={les $30$ instances de chaque régime, triées par delta croissant},
  ylabel={delta apparié des appels \ILP},
  xmin=0.5,xmax=30.5, ymin=-60, ymax=60,
  legend style={at={(0.98,0.03)},anchor=south east,font=\small,draw=cgris,
                fill=white,fill opacity=0.9,text opacity=1},
  grid=major, grid style={cgris!35}, tick label style={font=\small},
  label style={font=\small}]
\addplot[color=cdom,mark=*,mark size=1.4pt,thick]
  coordinates {(1,-205) (2,-81) (3,-15) (4,-4) (5,3) (6,3) (7,4) (8,5) (9,5) (10,6) (11,7) (12,7) (13,7) (14,8) (15,13) (16,14) (17,14) (18,16) (19,17) (20,17) (21,19) (22,21) (23,21) (24,21) (25,24) (26,25) (27,26) (28,27) (29,31) (30,51)};
\addlegendentry{$\text{corr}=0$ \ (front épais)}
\addplot[color=black!55,mark=triangle*,mark size=1.6pt,thick]
  coordinates {(1,-132) (2,-95) (3,-83) (4,-67) (5,-62) (6,-53) (7,-49) (8,-45) (9,-27) (10,-22) (11,-13) (12,-8) (13,-6) (14,-1) (15,1) (16,1) (17,1) (18,2) (19,7) (20,7) (21,8) (22,10) (23,10) (24,13) (25,13) (26,13) (27,15) (28,16) (29,21) (30,24)};
\addlegendentry{$\text{corr}=0{,}5$}
\addplot[color=ceff,mark=square*,mark size=1.3pt,thick]
  coordinates {(1,-129) (2,-48) (3,-36) (4,-32) (5,-21) (6,-15) (7,-12) (8,-12) (9,-10) (10,-9) (11,-9) (12,-9) (13,-9) (14,-6) (15,-6) (16,-6) (17,-3) (18,-3) (19,-3) (20,-3) (21,-1) (22,-1) (23,0) (24,0) (25,0) (26,2) (27,3) (28,4) (29,5) (30,8)};
\addlegendentry{$\text{corr}=0{,}9$ \ (front mince)}
\addplot[color=black,thick,dotted,domain=0.5:30.5] {0};
\end{axis}
\end{tikzpicture}
\caption{Lecture appariée du lot de $90$ instances : delta du nombre
d'appels au solveur entier, instance par instance, trié dans chaque régime.
Les valeurs hors cadre sont $-205$, $-132$, $-129$, $-95$, $-83$ et $-81$ --
toutes des \emph{baisses}. La figure dit ce qu'une médiane cache : à front
mince la courbe est presque entièrement sous zéro, à front épais presque
entièrement au-dessus, et le gain global tient à une queue courte de très
fortes baisses plutôt qu'à un progrès réparti. C'est pourquoi la médiane
appariée vaut $+1$ alors que le total baisse de $5{,}2\ \%$.}
\label{fig:apparie}
\end{figure}

\paragraph{Où le gain se trouve, et pourquoi.} Ventilé par corrélation entre
critères -- c'est-à-dire par épaisseur du front -- le signe est net et
conforme à la théorie.

\begin{table}[htbp]
\centering
\caption{Delta apparié des appels au solveur entier, par régime.}
\label{tab:corr}
\begin{tabular}{lccc}
\toprule
& $\text{corr}=0$ & $\text{corr}=0{,}5$ & $\text{corr}=0{,}9$\\
& (front épais) & & (front mince)\\
\midrule
médiane du delta      & $+14$ & $+1$ & $\mathbf{-6}$\\
total                 & $+107$ & $-501$ & $-361$\\
instances améliorées  & $4/30$ & $14/30$ & $\mathbf{22/30}$\\
\bottomrule
\end{tabular}
\end{table}

\noindent
Front mince : l'archive couvre vite $\Ecal$, le relâché se vide, et la
proposition~\ref{prop:vide} conclut à bas coût. Front épais : les points
dominés abondent, saturent le plafond, et chaque coupe d'archive posée
coûtait alors un programme entier de clôture pour rien -- c'est exactement
ce que le théorème~\ref{th:hauteur} supprime, en ramenant cette clôture à un
programme linéaire.

\subsection{À budget d'appels égal : un banc déterministe}
\label{sec:determ}

Les bancs précédents sont bornés en \emph{secondes}. Ils sont réalistes,
mais ils rendent invérifiable une affirmation telle que « aucune instance en
recul » : deux exécutions du même code n'y rendent pas le même nombre
d'appels. Le contrôle est sans appel -- même code, même graine, même
instance, trois exécutions à $18+12$~s :
\[
  q_{\text{lb}} = 26{,}9286\ ;\qquad 60{,}5000\ ;\qquad 60{,}5000 .
\]
Soit $125\ \%$ d'écart entre deux exécutions de la même chose.

Nous plafonnons donc en \textbf{nombre d'appels} au solveur entier, ce qui
retire l'horloge de la boucle de décision. Les mêmes trois exécutions
donnent alors trois fois la même valeur, le même nombre d'appels et la même
archive. Le banc \emph{vérifie} ce déterminisme au lieu de le supposer :
chaque configuration est jouée deux fois et toute ligne instable est exclue
de la conclusion.

\begin{table}[htbp]
\centering
\caption{Banc déterministe, plafond de $100$ appels, $90$ instances.}
\label{tab:determ}
\begin{tabular}{lc}
\toprule
reproductibilité & $\mathbf{90/90}$ exécutions identiques\\
preuves gagnées / perdues & $0$ / $\mathbf{0}$\\
valeurs plus faibles & $\mathbf{0}$\\
bornes plus lâches & $\mathbf{0}$\\
\bottomrule
\end{tabular}
\end{table}

\noindent
Deux conclusions, et la seconde est un résultat négatif qu'il faut énoncer
aussi nettement que les gains. La première : à budget d'appels égal, aucune
instance ne recule sur aucun des trois axes -- l'affirmation a désormais un
sens, et elle est établie. La seconde : \emph{aucune preuve n'est gagnée non
plus}, alors même que le plafond est serré ($6$ instances ne se prouvent
pas). Le théorème~\ref{th:hauteur} est donc bien ce que la
remarque~\ref{rem:lemme} annonçait : un gain de coût, et rien de plus.

\subsection{La coupe d'efficacité comme opérateur de recherche}
\label{sec:divers}

Le théorème~\ref{th:effcut} n'a pas servi qu'à borner. Couper autour de
\emph{tout} l'archive laisse exactement les points que l'archive ne domine
ni n'égale ; maximiser le substitut sur ce reste rend donc un point
\textbf{garanti neuf en vecteur critères} -- ce qu'aucun redémarrage
aléatoire ne garantit. On le répare en un point efficace certifié, on
l'archive, on coupe autour, on recommence.

\begin{remark}[aucune clôture n'est requise ici]
Employée comme opérateur de recherche, la coupe n'exige pas la
condition~\eqref{eq:cloture}. Celle-ci n'est nécessaire que pour tirer une
\emph{borne} d'un relâché amputé. Ici la coupe ne fait que diriger
l'exploration, et la validité du résultat ne tient qu'au test d'efficacité,
qui certifie chaque point rendu. Cette variante est donc \emph{gratuite} :
zéro programme entier de clôture.
\end{remark}

\begin{table}[htbp]
\centering
\caption{Diversification par coupes d'efficacité. $24$ paires,
$n=20$ à $50$, budget déterministe de $700$ appels, comparaison appariée.}
\label{tab:divers}
\begin{tabular}{lc}
\toprule
reproductibilité & $24/24$ lignes stables\\
instances \textbf{améliorées} & $\mathbf{6}$\\
instances dégradées & $\mathbf{0}$\\
inchangées & $18$\\
gain sur $q_{\text{lb}}$ & médiane $+0{,}00\ \%$, moyenne $+4{,}72\ \%$\\
                         & \textbf{maximum $+80{,}6\ \%$, minimum $+0{,}0\ \%$}\\
\bottomrule
\end{tabular}
\end{table}


\begin{figure}[htbp]
\centering
\begin{tikzpicture}
\begin{axis}[width=12.6cm,height=5.6cm,
  ybar, bar width=3.6pt,
  xlabel={les $24$ paires, triées par gain croissant},
  ylabel={gain sur $q_{\text{lb}}$ (\%)},
  xmin=0.3,xmax=24.7, ymin=-6, ymax=88,
  ytick={0,20,40,60,80},
  grid=major, grid style={cgris!35}, tick label style={font=\small},
  label style={font=\small}]
\addplot[fill=ceff,draw=ceff] coordinates {(1,0.0) (2,0.0) (3,0.0) (4,0.0) (5,0.0) (6,0.0) (7,0.0) (8,0.0) (9,0.0) (10,0.0) (11,0.0) (12,0.0) (13,0.0) (14,0.0) (15,0.0) (16,0.0) (17,0.0) (18,0.0) (19,2.7) (20,4.6) (21,4.9) (22,6.7) (23,13.8) (24,80.6)};
\addplot[color=black,thick,dotted,domain=0.3:24.7] {0};
\end{axis}
\end{tikzpicture}
\caption{Diversification par coupes d'efficacité, à budget d'appels
identique. Aucune barre ne descend sous zéro : sur les $24$ paires, la
valeur rendue n'est \emph{jamais} dégradée. Dix-huit paires sont
exactement nulles, six montent, de $+2{,}7$ à $+80{,}6\ \%$. C'est cette
asymétrie -- rien à perdre, parfois beaucoup à gagner -- qui justifie le
mécanisme, bien plus que sa moyenne.}
\label{fig:divers}
\end{figure}

\noindent
Le minimum vaut $+0{,}0\ \%$ : à budget d'appels égal, la diversification
n'a dégradé la valeur rendue sur \emph{aucune} des $24$ paires. Les six
gains vont de $+2{,}7$ à $+80{,}6\ \%$ et se répartissent sur $n=20$, $30$
et $40$ ainsi que sur les deux corrélations. Test des signes : $6$ contre
$0$, soit $p=0{,}016$.

Deux réserves, sans lesquelles ce tableau serait mal lu. D'abord il mesure
le \emph{mécanisme} et non son déclencheur : le réglage de production ne
l'active que si la borne paraît sans espoir, et sous ce réglage un banc
antérieur ne l'a jamais activé -- la sonde mesurait un écart de $23\ \%$ là
où le seuil est à $50\ \%$, et refusait donc de tirer, très correctement.
Ensuite il ne compare que $q_{\text{lb}}$ : les appels donnés à la
diversification sont autant d'appels retirés à la certification, et un
contrôle ponctuel montre que la borne \emph{supérieure}, elle, peut se
relâcher. Régler le déclencheur demande de mesurer les deux, ce qui reste à
faire.

\subsection{Où elle n'apporte rien, et pourquoi}

\begin{figure}[htbp]
\centering
\begin{tikzpicture}
\begin{axis}[width=12cm,height=5.4cm,
  xlabel={$n$}, ylabel={écart garanti (\%)},
  xmin=15,xmax=45, ymin=0, ymax=105,
  xtick={20,30,40}, ytick={0,20,40,60,80,100},
  legend style={at={(0.5,0.35)},anchor=north,font=\small,draw=cgris},
  grid=major, grid style={cgris!35}, tick label style={font=\small},
  label style={font=\small}]
\addplot[color=cdom,mark=*,thick,mark size=2.4pt]
  coordinates {(20,98.3) (30,97.1) (40,92.5)};
\addlegendentry{sans coupes d'archive}
\addplot[color=ceff,mark=square*,thick,dashed,mark size=2.2pt]
  coordinates {(20,98.3) (30,97.0) (40,92.4)};
\addlegendentry{avec}
\addplot[color=black,thick,dotted] coordinates {(15,0) (45,0)};
\addlegendentry{ce qu'il faudrait atteindre}
\end{axis}
\end{tikzpicture}
\caption{La frontière honnête de la méthode. À front épais et $n\ge 20$, les
deux courbes sont confondues et restent à $92$--$98\ \%$ : la coupe
d'efficacité n'y change \emph{rien}, et la borne ne descend pas. La raison
est mécanique -- la recherche produit assez de points dominés pour saturer
le plafond avant que l'archive ne soit atteinte, de sorte qu'aucune coupe
d'efficacité n'est posée. Nous reproduisons cette figure telle quelle plutôt
que de la reléguer à une note : elle délimite ce que ce travail
n'accomplit pas.}
\label{fig:echelle}
\end{figure}


À front épais et $n\ge 20$, les écarts garantis sont \textbf{inchangés} :
$98{,}3\ \%$, $97{,}1\ \%$ et $92{,}6\ \%$ à $n=20$, $30$, $40$. La raison
est mécanique et se lit dans la règle du vivier
(\S\ref{sec:vivier}) : à ces tailles la recherche produit assez de points
dominés pour saturer le plafond \emph{avant} que les points d'archive ne
soient atteints dans le classement. Aucune coupe d'efficacité n'est donc
posée.

C'est un résultat négatif utile, et il paraissait désigner le verrou suivant
sans ambiguïté : non plus la \emph{disponibilité} des coupes, mais le
\emph{plafond} lui-même, donc le coût de $p$ binaires par coupe. Désagréger
la disjonction~\eqref{eq:disj} devenait la priorité. La sous-section
suivante met cette conclusion à l'épreuve -- et l'infirme.

\subsection{Supprimer le coût en binaires ne suffit pas}
\label{sec:cglp}

\paragraph{Une coupe disjonctive sans aucune binaire.} Soit $P$ le relâché
courant dans l'espace des $x$ et, pour un point de coupe $a$, les $p$
disjoints
\[
  P_k=\bigl\{x\in P\ :\ \gamma_k^{\top}x\ \ge\ 1-\delta_k\bigr\},
  \qquad e^a_k(x)=\gamma_k^{\top}x+\delta_k .
\]
Tout point que l'on veut conserver appartient à leur \emph{réunion}. On
cherche donc $\alpha^{\top}x\ge\beta$ valide sur \emph{chacun} des $P_k$.
Par le lemme de Farkas, cela équivaut à l'existence de multiplicateurs
$\lambda^k\ge 0$ et $\mu^k\ge 0$ tels que
\[
  \alpha\ \ge\ -\lambda^k A+\mu^k\gamma_k,
  \qquad
  \beta\ \le\ -\lambda^k b+\mu^k(1-\delta_k),
\]
\emph{pour le même couple $(\alpha,\beta)$ à tous les $k$} -- c'est
précisément ce qui rend l'inégalité valide pour la réunion. On choisit alors
le couple le plus \emph{violé} au point courant $x^{\ast}$, ce qui donne un
programme linéaire de $n+1+p(m+1)$ variables. Un seul programme linéaire,
une seule ligne, \textbf{aucune binaire} : le plafond ne s'y applique pas.

\begin{remark}[ce qu'elle ne peut pas faire]
Elle ne domine pas la forme disjonctive. Celle-ci retranche \emph{exactement}
la région ; la coupe générée n'en retranche qu'un demi-espace. Son seul
avantage -- mais il était censé être décisif dans ce régime -- est de ne
consommer aucune place sous le plafond.
\end{remark}

\paragraph{Vérifiée avant d'être mesurée.} Sur dix instances énumérables :
aucune violation depuis le polyèdre de base, aucune en chaîne, et les seuls
points efficaces retranchés sont de même vecteur critères que leur point de
coupe -- la réserve du théorème~\ref{th:effcut}, non une faute. Un premier
contrôle avait signalé une violation ; elle n'en était pas une. Les coupes
s'enchaînent, et chacune est dérivée du polyèdre \emph{déjà} restreint par
les précédentes : elle n'est donc valide que relativement à \emph{ce}
polyèdre, et le point incriminé avait déjà été retiré par une coupe
antérieure. C'est une distinction qu'il faut faire, sous peine de rejeter
une implémentation correcte -- ou, plus grave, d'en accepter une fausse.

\begin{table}[htbp]
\centering
\caption{Coupe par programme générateur, budget déterministe de $300$ appels
entiers, comparaison appariée, chaque configuration jouée deux fois.}
\label{tab:cglp}
\setlength{\tabcolsep}{5pt}
\begin{tabular}{ccccccc}
\toprule
$n$ & corr & écart sans & écart avec & gain (points) & coupes générées
 & \textsc{Pl} sans/avec\\
\midrule
$20$ & $0$    & $98{,}6\ \%$ & $98{,}5\ \%$ & $+0{,}11$ & $13$ & $360/390$\\
$20$ & $0{,}5$& $0{,}0\ \%$  & $0{,}0\ \%$  & $+0{,}00$ & $0$  & $105/105$\\
$30$ & $0$    & $93{,}9\ \%$ & $93{,}8\ \%$ & $+0{,}06$ & $12$ & $315/346$\\
$30$ & $0{,}5$& $78{,}9\ \%$ & $78{,}9\ \%$ & $+0{,}00$ & $3$  & $270/279$\\
$40$ & $0$    & $96{,}6\ \%$ & $96{,}6\ \%$ & $+0{,}00$ & $12$ & $333/375$\\
$40$ & $0{,}5$& $0{,}0\ \%$  & $0{,}0\ \%$  & $+0{,}00$ & $9$  & $171/191$\\
\midrule
\multicolumn{7}{l}{reproductibilité $6/6$ \quad$\bullet$\quad
 améliorées $2$, dégradées $\mathbf{0}$, inchangées $4$}\\
\multicolumn{7}{l}{gain médian $+0{,}00$ point, moyen $+0{,}03$,
 \textbf{maximum $+0{,}11$}}\\
\bottomrule
\end{tabular}
\end{table}

\paragraph{Le mécanisme fonctionne, et cela ne suffit pas.} Les coupes sont
bien générées -- $13$, $12$, $12$, $9$, $3$ selon les lignes -- et elles
mordent : à $n=20$, $\text{corr}=0$, la borne supérieure passe de $230{,}26$
à $214{,}21$, soit $-7\ \%$. C'est la première fois qu'un levier déplace la
borne dans ce régime ; tout ce qui avait été essayé auparavant l'y laissait
strictement inchangée.

Mais l'écart garanti, lui, passe de $98{,}6$ à $98{,}5\ \%$. La raison est
arithmétique : avec $q_{\text{lb}}\approx 3{,}7$ et $q_{\text{ub}}\approx
230$, réduire la borne supérieure de $7\ \%$ ne déplace le rapport que d'un
dixième de point. Sur les six lignes, le gain médian est \emph{nul} et le
maximum vaut $0{,}11$ point, contre des écarts de $79$ à $99\ \%$.

\paragraph{Ce que ce résultat négatif corrige.} Il ne dit pas seulement que
la coupe générée est insuffisante ; il infirme le \emph{diagnostic} qui
l'avait motivée. Nous avions conclu de la section précédente que le verrou
était le coût en binaires. Or ce coût est ici entièrement supprimé -- les
coupes générées sont illimitées et n'occupent aucune place -- et l'écart ne
bouge pas. \textbf{Lever le plafond ne suffit donc pas}, et la conclusion
précédente était trop rapide.

Nous avions alors avancé une explication de remplacement : à ces tailles, la
région que l'on sait retrancher ne représenterait qu'une part négligeable de
ce qui sépare $\Rcal$ de $\Ecal$. La section suivante la met à l'épreuve --
et la réfute à son tour. Nous la laissons écrite ici, avec sa réfutation à la
page suivante, parce qu'une explication avancée puis abandonnée fait partie
du raisonnement et non de ses déchets.

% ============================================================================
% ============================================================================
\section{Ce que les coupes peuvent atteindre}
\label{sec:couverture}

Une première explication de l'écart résiduel a été avancée puis réfutée par
la mesure : le coût en binaires, donc le plafond de coupes
(§\ref{sec:cglp}). Une seconde restait : la région que l'on sait retrancher
serait négligeable devant ce qui sépare $\Rcal$ de $\Ecal$.
Cette section la met à l'épreuve. Elle demande de comparer ce que la méthode
retranche à ce qu'il \emph{faudrait} retrancher -- donc de connaître $\Ecal$
et $\qs$. Nous nous plaçons pour cela sur des instances énumérables, où
l'énumération exhaustive en arithmétique rationnelle exacte les fournit sans
la moindre approximation.

\subsection{La quantité à mesurer}

Posons
\[
  \Scal^{+}\ =\ \{\,x\in\Scal\ :\ f(x)>\qs\,\}.
\]
Ce sont exactement les points qui rendent la borne lâche : au-dessus de
$\qs$, une borne supérieure ne peut aller chercher qu'eux.

\begin{proposition}[l'outil de coupe est suffisant en principe]
\label{prop:suffisant}
Aucun point de $\Scal^{+}$ n'est efficace. Par conséquent, les coupes
d'efficacité centrées sur $\Ecal$ tout entier retranchent $\Scal^{+}$
en totalité.
\end{proposition}

\begin{proof}
Soit $x\in\Scal^{+}$. Si $x$ était efficace, on aurait $f(x)>\qs$ avec
$x\in\Ecal$, ce qui contredit $\qs=\max_{\Ecal}f$. Donc $x$ est dominé :
il existe $e\in\Ecal$ avec $Z(x)\le Z(e)$, $Z(x)\ne Z(e)$. Or la coupe
centrée en $e$ retranche exactement $\{y:Z(y)\le Z(e)\}$, qui contient $x$.
\end{proof}

L'outil n'est donc pas en cause : ce que l'on peut perdre, on le perd en
\emph{couverture}. Nous en mesurons deux niveaux,
\[
  \text{cov}_{\Acal}=\frac{|\{x\in\Scal^{+}:\exists a\in\Acal,\ Z(x)\le Z(a)\}|}
                         {|\Scal^{+}|},
  \qquad
  \text{cov}_{\text{coupes}}=\text{idem, centres réellement posés},
\]
le premier étant le plafond de l'archive, le second ce qui est effectivement
retranché. La lecture est sans ambiguïté : si $\text{cov}_{\Acal}$ s'effondre
avec $n$, le verrou est la couverture ; sinon, il est ailleurs.

\begin{remark}[arithmétique]
Les comparaisons de dominance sont exactes. Chaque coordonnée du vecteur
critères est remplacée par son \emph{rang} dans la liste triée des valeurs
distinctes prises sur $\Scal$ ; l'application étant strictement croissante
coordonnée par coordonnée, elle préserve l'ordre partiel. La comparaison
devient entière -- donc vectorisable -- sans rien perdre de l'exactitude des
fractions.
\end{remark}

\subsection{La mesure}

\begin{table}[htbp]
\centering
\caption{Couverture et bornes sur vérité terrain, $\text{corr}=0$, $p=3$,
plafond déterministe de $300$ appels entiers. $\kappa$ désigne l'échelle du
second membre. Échelle A à $\kappa$ fixé, le
seul régime rigoureusement comparable d'un $n$ à l'autre ; échelle B à
$\kappa$ décroissant pour garder $\Scal$ énumérable, indicative.}
\label{tab:couverture}
\setlength{\tabcolsep}{4pt}
\begin{tabular}{ccrrrccrrr}
\toprule
& $\kappa$ & $|\Scal|$ & $|Z(\Ecal)|$ & $|Z(\Acal)|$
 & $\text{cov}_{\Ecal}$ & $\text{cov}_{\Acal}$
 & $\qs$ & $q_{\text{ub}}$ & $U_{\text{coupes}}$\\
\midrule
\multicolumn{10}{l}{\emph{Échelle A -- $\kappa$ fixé}}\\
$6$  & $1{,}0$ & $250$    & $28$  & $16$ & $100\%$ & $100\%$ & $5{,}07$ & $5{,}07$ & $5{,}07$\\
$8$  & $1{,}0$ & $2\,387$ & $46$  & $15$ & $100\%$ & $100\%$ & $3{,}17$ & $3{,}17$ & $9{,}00$\\
$10$ & $1{,}0$ & $24\,152$& $70$  & $13$ & $100\%$ & $100\%$ & $37{,}67$& $37{,}67$& $39{,}33$\\
$12$ & $1{,}0$ & $304\,497$&$202$ & $24$ & $100\%$ & $100\%$ & $4{,}35$ & $\mathbf{32{,}27}$ & $\mathbf{8{,}53}$\\
\midrule
\multicolumn{10}{l}{\emph{Échelle B -- $\kappa$ décroissant}}\\
$12$ & $0{,}65$& $10\,486$& $78$  & $25$ & $100\%$ & $100\%$ & $4{,}15$ & $5{,}61$ & $4{,}15$\\
$14$ & $0{,}55$& $15\,411$& $127$ & $25$ & $100\%$ & $100\%$ & $5{,}00$ & $5{,}00$ & $5{,}75$\\
$16$ & $0{,}48$& $30\,861$& $94$  & $16$ & $100\%$ & $100\%$ & $11{,}50$& $11{,}50$& $28{,}50$\\
$18$ & $0{,}42$& $48\,703$& $139$ & $17$ & $100\%$ & $100\%$ & $14{,}20$& $14{,}20$& $17{,}60$\\
$20$ & $0{,}38$& $64\,268$& $136$ & $20$ & $100\%$ & $100\%$ & $8{,}13$ & $\mathbf{38{,}79}$ & $\mathbf{9{,}33}$\\
\bottomrule
\end{tabular}
\end{table}

Sur les $26$ lignes du lot complet ($n$ de $6$ à $20$, trois graines),
trois faits ressortent.

\paragraph{La proposition~\ref{prop:suffisant} se vérifie, $26$ fois sur
$26$.} $\text{cov}_{\Ecal}=100\ \%$ partout. Ce n'est pas une surprise --
c'est une démonstration -- mais c'est le contrôle qui donne son sens au
reste : si cette colonne avait montré autre chose, l'implémentation, et non
la théorie, aurait été en cause.

\paragraph{La couverture n'est pas le verrou, et l'archive est bien plus
efficace qu'on ne le croyait.} $\text{cov}_{\Acal}\ge 99\ \%$ sur $26$
lignes sur $26$, alors que l'archive ne retient qu'entre $20$ et $35\ \%$
des vecteurs critères de $\Ecal$ (médiane par échelle). Une petite archive
domine donc déjà la quasi-totalité de $\Scal^{+}$ : la seconde explication
tombe comme la première. Et le résultat a une portée propre, en
dehors du diagnostic : il dit qu'\emph{il ne sert à rien de chercher à
grossir l'archive}. Ce n'est pas en la remplissant qu'on gagnera.

\paragraph{Une perte que rien n'avait isolée : la recherche.} Sur $3$ des
$26$ lignes, toutes à $n\ge 12$, l'incumbent n'atteint pas $\qs$. Aucun banc
antérieur ne pouvait le voir : sans vérité terrain, un écart garanti de
$90\ \%$ ne dit pas s'il vient d'une borne trop haute ou d'un minorant trop
bas. Ici on lit les deux. À $n=12$, $\kappa=1$ : $\qs=4{,}35$ et
$q_{\text{lb}}=3{,}10$. Une partie de l'écart affiché n'est donc pas un
défaut de certification du tout.

\paragraph{Et le reste, que rien n'explique encore.} Regardons les deux
lignes en gras. À $n=12$, la borne annoncée vaut $32{,}27$ alors que le
maximum de $f$ sur la région \emph{réellement survivante} vaut $8{,}53$. À
$n=20$ : $38{,}79$ contre $9{,}33$. Les coupes ont fait leur travail --
$\text{cov}_{\text{coupes}}$ vaut $99{,}4\ \%$ et $100\ \%$ sur ces deux
lignes -- et la borne n'en tient pas compte. L'écart n'est donc ni dans le
nombre de coupes, ni dans leur placement, ni dans l'archive. Il est dans le
passage de la région à la borne. C'est là qu'il faut regarder, et c'est
l'objet de la section suivante.

% ============================================================================
% ============================================================================
\section{Le verrou était le seuil}
\label{sec:seuil}

La section précédente laisse une question précise : les coupes retranchent
la quasi-totalité de $\Scal^{+}$, et la borne n'en tient pas compte --
$32{,}27$ annoncé quand la région survivante plafonne à $8{,}53$. Le défaut
n'est donc pas dans le relâché mais dans la façon dont on en tire un nombre.

\subsection{Le théorème 5' ne parle pas de l'incumbent}

La certification pose $q=P/Q$, résout $U\ \ge\ \max_{\Rcal}F_q$ avec
$F_q=QN-PD$, et conclut $\qs\le q+U/(Q D^{+})$. Le seuil est \emph{toujours}
l'incumbent. Or rien dans la démonstration ne l'exige.

\begin{theorem}[borne à seuil libre]
\label{th:seuil}
Soit $t$ un rationnel quelconque et $G_t=N-tD$. Si l'on a
$\max_{x\in\Rcal}G_t(x)\le v$ pour un $v\ge 0$, alors
\[
  \max_{x\in\Rcal}f(x)\ \le\ t+\frac{v}{D_{\min}},
  \qquad D_{\min}=\min_{x\in\Scal}D(x)>0 .
\]
\end{theorem}

\begin{proof}
Pour $x\in\Rcal$ on a $N(x)-tD(x)\le v$, et $D(x)>0$ sous~(A1) ; donc
$f(x)=N(x)/D(x)\le t+v/D(x)$, et $v\ge0$ donne
$v/D(x)\le v/D_{\min}$. \end{proof}

Le théorème~5' est le cas $t=q$ : en effet $G_q=F_q/Q$, de sorte que
$v=U/Q$ et la borne s'écrit $q+U/(QD_{\min})$. Mais tout $t$ convient, et
\emph{chaque point du relâché en fournit un meilleur} : si l'ILP rend un
$x_t\in\Rcal$ avec $G_t(x_t)>0$, alors $f(x_t)>t$, et l'on recommence en
$t\leftarrow f(x_t)$. C'est un Dinkelbach ordinaire -- mené sur $\Rcal$ et
non sur $\Scal$. À la convergence ($v\le 0$), $\max_{\Rcal}f\le t$
exactement.

\begin{proposition}[validité]
\label{prop:geomvalide}
$\qs\ \le\ \max\bigl(q,\ \max_{\Rcal}f\bigr)$.
\end{proposition}

\begin{proof}
Les coupes de dominance préservent $\Ecal$ (Th.~4). Une coupe d'efficacité
n'est posée qu'après établissement de sa condition de clôture : les seuls
points efficaces qu'elle retranche ont le vecteur critères de son centre, et
aucun ne bat $q$. Donc tout $x\in\Ecal$ est soit dans $\Rcal$ -- et
$f(x)\le\max_{\Rcal}f$ -- soit retranché, et $f(x)\le q$. \end{proof}

\paragraph{Pourquoi la méthode ne le faisait pas.} La boucle de
certification n'avance son seuil que sur des points \emph{certifiés
efficaces}. Le réflexe est correct pour le minorant, où seul un point
efficace peut compter -- c'est la seule propriété que la matheuristique ne
peut pas perdre. Mais pour un \emph{majorant}, l'efficacité ne sert à rien :
n'importe quel point du relâché fait un seuil valide. Refuser d'avancer le
seuil sur un point non certifié coûtait donc tout l'écart entre $q$ et
$\max_{\Rcal}f$, sans rien acheter. C'est une garde utile appliquée au
mauvais endroit, et c'est le genre d'erreur qu'aucune relecture ne trouve :
il a fallu mesurer la région survivante pour la voir.

\begin{remark}[les deux routes sont incomparables]
La première utilise $D^{+}\ge D_{\min}$, plus fin, et la structure du test
d'efficacité -- elle peut donc descendre \emph{sous} $\max_{\Scal}f$, et
elle le fait. La seconde utilise un meilleur seuil. Aucune ne majore
l'autre : sur les $26$ lignes de vérité terrain, la seconde est meilleure
$3$ fois, moins bonne $20$ fois, égale $3$ fois -- mais ses trois victoires
sont précisément les lignes difficiles. Le minimum de deux bornes valides
étant valide, on les prend toutes les deux.
\end{remark}

\subsection{Ce que cela coûte}

La seconde route se paie sur le même plafond d'appels entiers : ce qu'on lui
donne, la boucle de coupes ne l'a pas. Deux tentatives de réglage ont été
faites, et le banc les a réfutées toutes les deux. Nous les rapportons parce
qu'elles disent quelque chose de la structure du problème.

\paragraph{Première tentative : décider après coup.} Réserver \emph{en
aveugle} paraissait un pari ; nous avons voulu prendre la décision au moment
où l'information existe, c'est-à-dire après le premier tour, quand la
première route a produit une borne. Cela ne peut pas fonctionner ici : le
premier tour consomme à lui seul la \emph{totalité} du plafond -- l'oracle
interne tourne jusqu'à épuisement -- de sorte qu'à sa fin il ne reste rien à
réserver. Le symptôme était sans équivoque : statut « interrompue », zéro
tour, zéro appel, gain $+0{,}00$ là où la réserve en amont gagnait $9{,}82$
points. Le mécanisme ne choisissait pas mal, il ne s'exécutait pas.

\paragraph{Seconde tentative : réserver puis libérer.} La réserve est donc
prise avant la boucle, et le déclencheur change de sens -- il ne réserve
plus, il \emph{libère} : si le premier tour montre une borne déjà serrée, la
part réservée est rendue et le tour suivant en profite. Le banc l'a réfuté
aussi, et pour la même raison de séquence : quand la libération se décide,
le premier tour a \emph{déjà} tourné sous le plafond réduit, et rendre la
part ne rend pas les coupes qu'il n'a pas posées.

\begin{table}[htbp]
\centering
\caption{Le déclencheur contre son absence, mêmes douze instances, même
plafond de $300$ appels entiers, comparaison appariée.}
\label{tab:porte}
\begin{tabular}{lcccc}
\toprule
& améliorées & dégradées & inchangées & gain médian\\
\midrule
\textbf{sans déclencheur} & $\mathbf{8}$ & $\mathbf{1}$ & $3$
 & $\mathbf{+3{,}19}$\\
avec déclencheur & $6$ & $2$ & $4$ & $+0{,}03$\\
\bottomrule
\end{tabular}
\end{table}

Le détail est plus net que le résumé. À $n=30$, $\text{corr}=0$, graine~$2$,
le déclencheur refuse d'engager la route -- et la ligne perd tout de même
$0{,}26$ point, contre un \emph{gain} de $6{,}04$ sans lui. Il paie le coût
et refuse le bénéfice. À $n=20$, $\text{corr}=0$, graine~$2$, il abandonne
$+0{,}33$. Il est donc désactivé par défaut ; le code en est conservé,
puisqu'un mécanisme dont on a mesuré qu'il nuit se documente mieux en place
qu'effacé.

\paragraph{Ce qui a réellement corrigé le troc.} Non pas un déclencheur mais
la \emph{taille} de la réserve. La première version en réservait un
pourcentage -- $25\ \%$ du budget restant, soit $28$ appels -- à une
procédure qui en dépense $2$ à $4$. Vingt-six appels étaient retirés aux
coupes sans contrepartie. Bornée par le nombre de tours que la route peut
effectivement faire, la réserve ne coûte plus que $3$ à $4$ coupes. Le
paramètre n'a pas été réglé sur le banc : il a été \emph{dérivé} de la
procédure qu'il finance.

\subsection{Ce que cela déplace}

\begin{table}[htbp]
\centering
\caption{Seconde route, budget déterministe de $300$ appels entiers,
comparaison appariée, chaque configuration jouée deux fois. « tours » et
« \ILP{}~2\textsuperscript{e} » donnent le coût \emph{propre} de la seconde
route ; « coupes » ce que la boucle a cédé pour la financer.}
\label{tab:geom}
\setlength{\tabcolsep}{4pt}
\begin{tabular}{ccccrccc}
\toprule
$n$ & corr & écart sans & écart avec & gain (pts) & tours
 & coupes s/a & \ILP{} s/a\\
\midrule
$20$ & $0$     & $98{,}6\ \%$ & $88{,}8\ \%$ & $\mathbf{+9{,}82}$  & $2$ & $119/115$ & $300/290$\\
$20$ & $0$     & $42{,}7\ \%$ & $42{,}3\ \%$ & $+0{,}33$           & $2$ & $108/105$ & $300/290$\\
$20$ & $0{,}5$ & $0{,}0\ \%$  & $0{,}0\ \%$  & $+0{,}00$           & --  & $32/32$   & $143/143$\\
$20$ & $0{,}5$ & $0{,}0\ \%$  & $0{,}0\ \%$  & $+0{,}00$           & --  & $57/57$   & $217/217$\\
$30$ & $0$     & $93{,}9\ \%$ & $78{,}6\ \%$ & $\mathbf{+15{,}26}$ & $4$ & $104/101$ & $300/292$\\
$30$ & $0$     & $33{,}2\ \%$ & $27{,}1\ \%$ & $+6{,}04$           & $2$ & $110/107$ & $300/290$\\
$30$ & $0{,}5$ & $78{,}9\ \%$ & $\mathbf{8{,}0\ \%}$ & $\mathbf{+70{,}86}$ & $3$ & $89/85$ & $300/291$\\
$30$ & $0{,}5$ & $12{,}7\ \%$ & $13{,}7\ \%$ & $-1{,}00$           & $2$ & $97/93$   & $300/290$\\
$40$ & $0$     & $96{,}6\ \%$ & $96{,}5\ \%$ & $+0{,}06$           & $3$ & $110/107$ & $300/291$\\
$40$ & $0$     & $94{,}0\ \%$ & $87{,}3\ \%$ & $+6{,}74$           & $4$ & $98/95$   & $300/292$\\
$40$ & $0{,}5$ & $0{,}0\ \%$  & $0{,}0\ \%$  & $+0{,}00$           & --  & $56/56$   & $197/197$\\
$40$ & $0{,}5$ & $82{,}6\ \%$ & $43{,}9\ \%$ & $\mathbf{+38{,}64}$ & $4$ & $93/90$   & $300/292$\\
\midrule
\multicolumn{8}{l}{reproductibilité $12/12$ \quad$\bullet$\quad
 cohérence des encadrements : toutes valides}\\
\multicolumn{8}{l}{améliorées $\mathbf{8}$, dégradées $\mathbf{1}$,
 inchangées $3$ \quad$\bullet$\quad convergée $9/12$
 \quad$\bullet$\quad test des signes $p=0{,}039$}\\
\multicolumn{8}{l}{gain médian $\mathbf{+3{,}19}$ point, moyen
 $\mathbf{+12{,}23}$, maximum $+70{,}86$, minimum $-1{,}00$}\\
\bottomrule
\end{tabular}
\end{table}

\paragraph{L'ordre de grandeur change.} Le levier précédent -- la coupe
générée sans binaire -- rapportait au mieux $0{,}11$ point
(§\ref{sec:cglp}). Celui-ci en rapporte $70{,}86$ sur sa meilleure ligne,
$38{,}64$ sur la suivante, $15{,}26$ et $9{,}82$ ensuite. Ce n'est pas le
même phénomène mesuré plus finement, c'est un autre phénomène. Et il agit là
où la section d'échelle disait que rien n'agissait : à $n=30$,
$\text{corr}=0{,}5$, l'écart garanti tombe de $78{,}9$ à $8{,}0\ \%$ --
c'est-à-dire qu'on passe d'une borne qui n'apprenait rien à un encadrement
serré. Sur les $9$ lignes non nulles, $8$ progressent et une seule recule :
le test des signes donne $p=0{,}039$.

\paragraph{Elle coûte deux à quatre appels entiers.} C'est la colonne
« tours », et c'est le chiffre à retenir : un tour de Dinkelbach est un ILP,
et la route converge en $2$ à $4$ tours. Le reste de l'écart entre $300$ et
$290$ appels est sa \emph{réserve}, non sa consommation.

\paragraph{Une ligne recule encore.} $-1{,}00$ point, sur une instance dont
l'écart valait déjà $12{,}7\ \%$ : la première route y fermait convenablement
et les trois coupes cédées lui manquent. C'est le résidu du même arbitrage,
réduit d'un facteur deux par la réserve bornée, et il ne disparaît pas
complètement.

\paragraph{Une limite honnête, et la question qu'elle ouvre.} À $n=40$,
$\text{corr}=0$, graine~$1$, l'écart ne bouge que de $0{,}06$ point --
$96{,}6$ contre $96{,}5$ -- alors même que le statut « convergée » certifie
que la borne annoncée \emph{est} $\max_{\Rcal}f$. Si le majorant ne peut
plus descendre, ce qui reste tient à l'autre terme de l'écart :
$q_{\text{lb}}$. Or la section~\ref{sec:couverture} a précisément établi, là
où la vérité terrain existe, que l'incumbent manque $\qs$ sur $3$ lignes de
$26$, toutes à $n\ge 12$. L'hypothèse est donc que ce qui reste à ces
tailles n'est plus un défaut de certification mais un défaut de
\emph{recherche} -- ce qui déplacerait l'effort à l'opposé exact de là où
nous l'avons porté jusqu'ici. Nous ne l'avons pas encore mesuré.

% ============================================================================
% ============================================================================
\section{Comparaison avec les méthodes du domaine}
\label{sec:comparaison}

\subsection{Trois problèmes voisins, et pourquoi aucun n'est le nôtre}

Deux méthodes exactes encadrent directement ce travail. Il faut commencer
par dire, très précisément, ce qu'elles résolvent -- car cela commande la
forme que toute comparaison peut prendre.

\begin{table}[htbp]
\centering
\caption{Le problème traité, chez les uns et chez les autres.}
\label{tab:positionnement}
\begin{tabular}{lcc}
\toprule
& critères $\Zk$ & fonction à optimiser sur $\Ecal$\\
\midrule
Zerdani \& Moulaï (2011) & \textbf{fractionnaires} & linéaire\\
Drici, Ouail \& Moulaï (2018) & linéaires & \textbf{fractionnaire}\\
\midrule
\textbf{notre problème (P)} & \textbf{fractionnaires} & \textbf{fractionnaire}\\
\bottomrule
\end{tabular}
\end{table}

\noindent
Notre problème est la \emph{réunion} des deux généralisations : fractionnaire
des deux côtés. Aucune des deux méthodes ne le couvre. Il en découle une
contrainte méthodologique dont ce mémoire tire parti plutôt que de la
subir : \textbf{la comparaison ne peut se faire que sur le terrain de
chacun}. On restreint donc notre méthode à leur classe -- $f$ linéaire pour
l'une, critères linéaires pour l'autre -- et on les affronte là où elles
sont chez elles. C'est la forme la plus exigeante de comparaison : si notre
méthode tient dans la spécialité de l'autre, le cas général suit
\emph{a fortiori}.

\subsection{Validation externe : reproduire, avant de comparer}
\label{sec:externe}

Avant de mesurer quoi que ce soit, il faut établir que nos objets sont bien
les leurs. Les deux articles donnent un exemple entièrement résolu ; nous
les reprenons à l'identique.

\begin{table}[htbp]
\centering
\caption{Les deux exemples publiés, reproduits.}
\label{tab:exemples}
\begin{tabular}{lcc}
\toprule
& Zerdani \& Moulaï, \S5 & Drici \emph{et al.}, \S5\\
\midrule
$|\Scal|$ / $|\Ecal|$ & $6$ / $5$ & $14$ / $7$\\
ensemble efficace & \textbf{identique au publié} & les deux optima publiés\\
                  &                              & sont efficaces chez nous\\
optimum publié & $\varphi=6$ en $(3,0)$ & $f=-5$\\
optimum calculé & $\varphi=6$ en $(3,0)$ & $f=-5$\\
notre méthode & prouve $6$, $10$ \ILP & prouve $-5$, $101$ \ILP\\
\bottomrule
\end{tabular}
\end{table}

\noindent
Un contrôle supplémentaire confirme la transcription de la seconde instance :
les auteurs donnent l'optimum de la relaxation continue, $x=(0,\frac{14}{5},
\frac{2}{5},0)$ de valeur $-\frac{85}{21}$ ; notre instance rend exactement
cette valeur en ce point.

\paragraph{Ce que la reproduction a révélé, et qui n'était pas prévu.} La
première exécution sur l'exemple de Drici \emph{et al.} donnait $|\Scal|=12$
et un optimum de $-\frac{19}{4}$, en désaccord avec les $-5$ publiés ; pis,
l'un des deux optima publiés, $(2,3,0,0)$, n'apparaissait même pas efficace.
Avant de mettre en cause l'article, nous avons compté $\Scal$ à la main :
$14$ points, pas $12$. Les manquants étaient tous ceux à $x_1=2$, dont
$(2,3,0,0)$ lui-même.

La faute était dans \emph{notre} énumérateur. Son élagage coupait le
sous-arbre dès qu'une somme partielle dépassait $b$ -- raisonnement exact
sous (A2), que notre générateur garantit, si bien que rien ne s'était jamais
vu sur les instances tirées. Mais les instances de la littérature ont des
coefficients \emph{négatifs} : une variable ultérieure peut alors faire
redescendre une ligne sous $b$, et l'élagage retirait des points réalisables.

La portée mérite d'être énoncée : \texttt{ground\_truth} est la référence de
\emph{toutes} les vérifications de ce travail, et elle était fausse
exactement là où l'on validait contre l'extérieur. L'élagage a été remplacé
par une minoration valide de la contribution des variables encore libres ;
les campagnes sur instances générées sont inchangées, ce que confirme la
batterie $V_1$--$V_4$.

\subsection{Réimplémenter, plutôt que citer}
\label{sec:reimpl}

Comparer à une méthode exige de la faire tourner. Celle de Zerdani \& Moulaï
ne se lit pas dans la valeur d'un optimum mais dans le \textbf{tableau du
simplexe} : l'ensemble hors base $N_k$, le gradient réduit $\gamma_{k,j}$,
les colonnes $y_{k,ij}$, les arêtes $E_{j_k}$ qu'elles définissent, le pas
entier $\theta_{j_k}$ le long de ces arêtes. Aucun solveur industriel
n'expose cela.

\paragraph{Un simplexe rationnel, et pourquoi rationnel.} Les tests de leur
méthode portent sur des \emph{égalités} : « $\gamma_{k,j}=0$ » décide de
l'unicité de l'optimum -- donc de son efficacité par leur corollaire~2.3 --
et « $\theta$ entier » décide de l'existence d'un point entier sur une
arête. En flottant, ces tests réclament un $\varepsilon$, et un $\varepsilon$
mal choisi change le \emph{résultat} de l'algorithme, non sa précision. Les
données étant entières, tout le tableau reste rationnel : l'exactitude
s'obtient sans rien supposer. L'objectif fractionnaire est traité par le
gradient réduit de Martos, sous la forme même de leur section~2 ; la règle
de Bland des deux côtés garantit la terminaison sans perturbation, la
dégénérescence étant fréquente sur ces données. Une phase~I a dû être
ajoutée : leurs coupes $\sum_{j\in N_k}x_j\ge 1$ rendent la base des écarts
irréalisable, ce pour quoi ils prescrivent le simplexe dual.

\paragraph{Deux endroits où il a fallu trancher.} Le premier est le test
d'efficacité : les auteurs invoquent leur théorème~3.6, nous employons notre
test intégral, qui décide le \emph{même} prédicat de façon exacte. Le second
est le sommet fractionnaire : les étapes 3.2 et 3.3 prescrivent « dual
simplex and Gomory cuts, if necessary » sans détailler ces coupes. Nous
avons implémenté la coupe fractionnaire de Gomory en arithmétique exacte --
c'est bien ce qu'ils prescrivent -- et signalé les cas où le processus
n'aboutit pas malgré elle.

\paragraph{Un bogue de portée, et ce qu'il a coûté.} Une fois les coupes de
Gomory en place, notre réimplémentation rendait $\varphi_{\text{opt}}=8$ en
$(4,0)$ \emph{sur leur propre exemple}, contre les $6$ en $(3,0)$ publiés --
et $(4,0)$ est réalisable mais \emph{dominé}. La cause : nous appliquions le
corollaire~2.3 à chaque itération, alors qu'il porte sur le domaine
d'\emph{origine}. Dès qu'une coupe est posée, le sommet n'optimise plus
$Z_1$ que sur la région tronquée et le corollaire ne s'applique plus ; les
auteurs ne s'en servent d'ailleurs qu'à l'étape~2, et testent explicitement
l'efficacité à l'étape~3.3.

C'est le genre d'erreur qui rend une réimplémentation inutilisable sans
qu'on s'en aperçoive : elle tournait, terminait, et rendait un nombre
plausible. Seule la confrontation à l'exemple publié l'a révélée -- ce qui
justifie rétrospectivement d'avoir commencé par là. Après correction :
statut \emph{optimal}, $\varphi_{\text{opt}}=6$ en $(3,0)$, en $7$
itérations, $6$ coupes et $14$ tests d'efficacité.

\subsection{Le match sur le terrain de Zerdani \& Moulaï}
\label{sec:matchzm}

\paragraph{Protocole.} $24$ instances de leur classe -- $\varphi$ linéaire,
critères fractionnaires -- avec vérité terrain par énumération exhaustive.
L'unité de coût est le nombre d'appels au solveur \emph{en nombres entiers},
indépendante de la machine et du langage : les secondes ne compareraient
rien, un simplexe rationnel exact d'un côté et HiGHS de l'autre. Le plafond
de travail de leur méthode porte sur le nombre d'itérations et de coupes
($60$ et $150$), jamais sur des secondes, pour rester reproductible.

\begin{remark}[l'unité choisie les favorise, et c'est voulu]
Leur algorithme tire l'essentiel de son travail du tableau -- unicité lue
dans $\gamma$, arêtes lues dans les colonnes -- opérations que cette unité
ne compte \emph{pas}, alors qu'elle compte tout ce que nous faisons. Le
désavantage est assumé : si notre méthode tient malgré cela, la conclusion
est solide.
\end{remark}

\begin{table}[htbp]
\centering
\caption{Match, $24$ instances, vérité terrain par énumération.}
\label{tab:match}
\begin{tabular}{lcc}
\toprule
& Zerdani \& Moulaï & notre méthode\\
\midrule
optimum \emph{trouvé} & $20/24$ & $\mathbf{24/24}$\\
optimalité \textbf{prouvée} & $12/24$ & $\mathbf{24/24}$\\
\quad dont trouvé mais non certifié & $8$ & --\\
coût médian & $45$ tests d'efficacité & $79$ \ILP\\
encadrement valide & \multicolumn{2}{c}{$24/24$}\\
\bottomrule
\end{tabular}
\end{table}

\paragraph{Premier enseignement : un raccourci qui leur appartient.} Sur
$9$ instances sur $24$, leur méthode conclut en \textbf{zéro coupe et un
seul test d'efficacité}. L'optimum de $\varphi$ sur $\Scal$, calculé à
l'étape~0, se trouve déjà efficace, et l'étape~1 termine. Nous dépensons
entre $25$ et $158$ appels sur ces mêmes instances.

C'est un avantage réel de leur construction, et il faut le dire nettement :
sur ces instances-là, ils sont imbattables. Il tient à ce que $\varphi$ est
\emph{linéaire} -- un seul programme entier donne son maximum sur $\Scal$,
et il n'y a plus qu'à tester ce point. Notre $f$ fractionnaire n'a pas
d'équivalent de cette étape : le maximum de $f$ sur $\Scal$ demande déjà un
Dinkelbach complet. Aucune de nos mesures antérieures ne pouvait révéler cet
écart, puisqu'elles ne comparaient la méthode qu'à elle-même.

\paragraph{Deuxième enseignement : une coupure nette selon $|\Scal|$.}
Quand leur processus s'achève, $|\Scal|$ médian vaut $121$ ; quand il ne
s'achève pas, $371$. Le mécanisme est identifié et tient à la coupe de
Dantzig : elle retire \textbf{un sommet à la fois}. La certification
parcourt donc les sommets, et son coût suit $|\Scal|$ plutôt que la
difficulté du problème. On observe le phénomène directement sur une instance
à $|\Scal|=329$ : l'optimum $\varphi^\star=15$ est atteint dès la
\emph{première} itération, et les $150$ coupes suivantes ne servent qu'à
tenter de le certifier.


\begin{figure}[htbp]
\centering
\begin{tikzpicture}
\begin{semilogxaxis}[width=12.6cm,height=5.4cm,
  xlabel={$|\Scal|$ (échelle logarithmique)},
  xmin=15,xmax=6000, ymin=-0.6, ymax=2.6,
  ytick={0,1,2},
  yticklabels={{n'atteint pas $\varphi^\star$},{$\varphi^\star$ atteint,
    non certifié},{processus achevé}},
  yticklabel style={font=\small,align=right},
  grid=major, grid style={cgris!35}, tick label style={font=\small},
  label style={font=\small}]
\addplot[only marks,mark=*,mark size=2.4pt,color=ceff]
  coordinates {(23,2) (94,2) (22,2) (412,2) (19,2) (1450,2) (32,2) (121,2) (230,2) (24,2) (185,2) (1428,2)};
\addplot[only marks,mark=square*,mark size=2.2pt,color=cdom]
  coordinates {(329,1) (77,1) (371,1) (4657,1) (25,1) (81,1) (527,1) (341,1)};
\addplot[only marks,mark=triangle*,mark size=2.8pt,color=black]
  coordinates {(1914,0) (1896,0) (124,0) (2116,0)};
\end{semilogxaxis}
\end{tikzpicture}
\caption{Match sur le terrain de Zerdani \& Moulaï : ce que rend leur
méthode, en fonction de la taille du domaine entier. La séparation est
nette. Quand le processus s'achève (en haut), $|\Scal|$ vaut $121$ en
médiane ; quand il n'y parvient pas, $371$. La cause est mécanique : la
coupe de Dantzig retire \emph{un sommet à la fois}, de sorte que le coût de
la certification suit $|\Scal|$ et non la difficulté du problème. Notre
méthode, elle, se situerait entièrement sur la ligne du haut : elle prouve
les $24$.}
\label{fig:matchzm}
\end{figure}

\paragraph{Troisième enseignement : ce que rend une méthode exacte
interrompue.} Sur $4$ instances, leur méthode ne trouve pas l'optimum dans
le plafond ($35$ au lieu de $39$ ; $-7$ au lieu de $8$ ; $-17$ au lieu de
$-12$ ; $1$ au lieu de $8$), toutes à $|\Scal|$ grand. Sur $8$ autres elle
le trouve sans pouvoir le certifier. Au total, sur \textbf{$12$ instances
sur $24$}, elle ne rend rien d'exploitable au plafond -- non par défaut de
mise en \oe uvre, mais par nature : une méthode exacte rend l'optimum
certifié ou rien.

C'est exactement l'axe sur lequel les deux approches diffèrent, et le seul
qui ne dépende ni de la machine ni de l'implémentation. Notre méthode rend
en permanence un couple $(q_{\text{lb}}, q_{\text{ub}})$ : sur ces mêmes
$12$ instances elle rend une solution \emph{et} un écart garanti, ici nul
puisqu'elle les prouve toutes.

\subsection{Le terrain de Drici, Ouail \& Moulaï}
\label{sec:matchdrici}

Sur l'autre bord -- critères linéaires, $f$ fractionnaire -- nous reprenons
leur protocole de génération, cité à leur section~6 : contraintes uniformes
sur $[1,30]$, second membre sur $[50,100]$, coefficients des critères et du
numérateur de $f$ sur $[-10,10]$, et $q$, $\beta$ tels que
$q^{\top}x+\beta>0$ sur le domaine. Leurs tailles sont reprises telles
quelles.

\begin{table}[htbp]
\centering
\caption{Notre méthode sur le protocole et les tailles de Drici
\emph{et al.}, $10$ instances par taille, budget déterministe de $1500$
appels.}
\label{tab:drici}
\begin{tabular}{lcccc}
\toprule
$n\times m$ & optimalité prouvée & écart garanti médian & \ILP\ méd.
 & CPU publié\\
\midrule
$5\times 5$   & $10/10$ & $0{,}0\ \%$ & $97$  & $0{,}21$~s\\
$10\times 5$  & $10/10$ & $0{,}0\ \%$ & $219$ & $0{,}44$~s\\
$15\times 5$  & $10/10$ & $0{,}0\ \%$ & $166$ & $0{,}70$~s\\
$20\times 5$  & $10/10$ & $0{,}0\ \%$ & $156$ & $1{,}23$~s\\
$20\times 10$ & $10/10$ & $0{,}0\ \%$ & $216$ & $1{,}45$~s\\
$30\times 10$ & $10/10$ & $0{,}0\ \%$ & $340$ & $8{,}98$~s\\
$35\times 15$ & $10/10$ & $0{,}0\ \%$ & $376$ & $7{,}45$~s\\
$40\times 15$ & $10/10$ & $0{,}0\ \%$ & $352$ & $10{,}82$~s\\
\bottomrule
\end{tabular}
\end{table}

\noindent
Quatre-vingts instances sur quatre-vingts à optimalité prouvée, et
$q_{\text{lb}}\le\qs\le q_{\text{ub}}$ partout. La colonne « CPU publié »
est le temps moyen rapporté par les auteurs (\textsc{Matlab}, MacBook Pro,
$r=3$) : elle situe l'ordre de grandeur des tailles atteintes, elle
\textbf{ne départage pas} des implémentations mesurées sur des machines et
des langages différents à des années d'intervalle, et nous ne l'utilisons
pas ainsi.

\paragraph{Une réserve qui doit précéder toute lecture de ce tableau.} Le
couple $(n,m)$ ne mesure pas la difficulté : les \emph{bornes des variables}
la mesurent. Des contraintes sur $[1,30]$ avec un second membre sur
$[50,100]$ donnent des variables bornées par $2$ ou $3$.

\begin{table}[htbp]
\centering
\caption{« $n=40$ » ne désigne pas le même problème d'un protocole à
l'autre.}
\label{tab:boite}
\begin{tabular}{lcccc}
\toprule
protocole & $n$ & $\overline{x}$ max & $\overline{x}$ méd.
 & $\log_{10}$ de la boîte\\
\midrule
Drici \emph{et al.} ($A\sim[1,30]$) & $40$ & $3$ & $2$ & $18{,}2$\\
notre lot d'échelle & $40$ & $24$ & $19$ & $\mathbf{51{,}8}$\\
\bottomrule
\end{tabular}
\end{table}


\begin{figure}[htbp]
\centering
\begin{tikzpicture}
\begin{axis}[width=11.4cm,height=5.2cm,
  ybar, bar width=13pt,
  symbolic x coords={$n=20$,$n=30$,$n=40$},
  xtick=data,
  ylabel={$\log_{10}$ du nombre de points de la boîte},
  ymin=0, ymax=58,
  nodes near coords, nodes near coords style={font=\scriptsize},
  legend style={at={(0.03,0.97)},anchor=north west,font=\small,draw=cgris},
  grid=major, grid style={cgris!35}, tick label style={font=\small},
  label style={font=\small}]
\addplot[fill=cdom!45,draw=cdom] coordinates
  {($n=20$,7.8) ($n=30$,13.5) ($n=40$,18.2)};
\addlegendentry{protocole de Drici \emph{et al.}}
\addplot[fill=ceff!35,draw=ceff] coordinates
  {($n=20$,20.9) ($n=30$,36.2) ($n=40$,51.8)};
\addlegendentry{notre lot d'échelle}
\end{axis}
\end{tikzpicture}
\caption{Pourquoi « $n=40$ » ne désigne pas le même problème. Des
contraintes tirées sur $[1,30]$ avec un second membre sur $[50,100]$
bornent les variables par $2$ ou $3$ ; les nôtres le sont par $19$ en
médiane. À la même valeur de $n$, trente-quatre ordres de grandeur séparent
les deux boîtes. Prouver $80/80$ sur le protocole de gauche ne dit rien du
régime de droite -- et c'est ce dernier qui nous résiste.}
\label{fig:boite}
\end{figure}

\noindent
Trente-quatre ordres de grandeur séparent les deux boîtes à la même valeur
de $n$. Prouver $80/80$ ici ne contredit donc \emph{en rien} les
$92\ \%$ d'écart garanti à $n=40$ mesurés à la
section~\ref{sec:res} sur nos propres instances, et il serait malhonnête de
présenter ce tableau comme une résolution du régime qui nous résiste. Les
deux mesures portent sur des régimes différents et se lisent séparément.

\subsection{Ce que la comparaison établit, et ce qu'elle n'établit pas}

\textbf{Établi.} Sur la classe de Drici \emph{et al.}, avec leur règle de
génération et leurs tailles, la méthode prouve l'optimalité partout. Sur
celle de Zerdani \& Moulaï, elle prouve $24/24$ là où leur méthode certifie
$12/24$, et cela sous une unité de coût qui leur est favorable. Les deux
exemples publiés sont reproduits exactement.

\textbf{Non établi.} Rien n'est dit des temps de calcul, qui ne se comparent
pas ici. Rien n'est dit non plus des grandes instances \emph{de notre}
protocole, où aucune des trois méthodes ne conclut. Et le raccourci de
l'étape~1 de Zerdani \& Moulaï reste un avantage net de leur construction
sur les instances où il s'applique : notre fonction objectif étant
fractionnaire, nous n'en disposons pas.

\section{Conclusion}

La coupe d'efficacité repose sur une observation d'une ligne -- deux points
efficaces de vecteurs critères distincts se dépassent forcément quelque part
-- et sur le fait, non trivial, que la réserve qui l'accompagne se vérifie
par \emph{un seul programme de faisabilité}. Elle transforme un sous-produit
de la recherche, l'archive, en la ressource dont la borne exacte manquait.

C'est une hybridation au sens fort, et il faut le dire précisément :
\emph{aucune des deux moitiés ne produit seule ce résultat}. Sans la
recherche heuristique il n'y a pas d'archive. Sans la certification exacte
de l'efficacité, les points de l'archive ne seraient pas \emph{prouvés}
efficaces et la coupe serait invalide. La proposition~\ref{prop:vide} est la
signature de cette hybridation : une preuve d'optimalité que la partie exacte
seule ne peut pas atteindre, puisqu'elle repose sur le fait d'avoir
\emph{épuisé} $\Ecal$ par des points trouvés heuristiquement.

\subsection*{Ce que la méthode ne fait pas}

Trois limites, énoncées avec la même précision que les gains, parce qu'elles
délimitent exactement ce que ce travail établit.

\begin{enumerate}[leftmargin=*,itemsep=3pt]
\item \textbf{Aucun effet au-delà de $n\ge 20$.} Les écarts garantis restent
      à $98{,}3$, $97{,}1$ et $92{,}6\ \%$, inchangés. Dans ce régime la
      recherche produit assez de points dominés pour saturer le plafond
      avant que l'archive ne soit atteinte : aucune coupe d'efficacité n'est
      posée. Le verrou n'est plus la \emph{disponibilité} des coupes mais
      leur \emph{coût}, $p$ binaires chacune -- et la
      section~\ref{sec:cglp} montre que lever ce coût ne suffit pas
      davantage.
\item \textbf{Aucun progrès sur la qualité des solutions.} Sur le lot figé
      la recherche atteint déjà l'optimum $90$ fois sur $90$ : il n'y a pas
      de place pour un tel progrès, et la contribution porte donc
      entièrement sur la \emph{certification}. Aux tailles où de la place
      existerait, nous ne disposons pas encore d'une mesure valide.
\item \textbf{Le théorème~\ref{th:hauteur} n'achète pas de preuves.} À
      budget d'appels égal il n'en gagne aucune (\S\ref{sec:determ}) : c'est
      un gain de coût, non de capacité.
\end{enumerate}

\subsection*{Directions ouvertes}

La première n'est plus celle que nous pensions. Nous avions désigné le
plafond de coupes -- donc le coût de $p$ binaires -- comme le verrou
suivant. La section~\ref{sec:cglp} lève ce coût entièrement, par un
programme générateur produisant des coupes disjonctives valides sans aucune
binaire, en nombre illimité : l'écart garanti gagne $0{,}11$ point au mieux,
zéro en médiane. Le plafond n'était donc pas le verrou.

Ce que la mesure met à sa place est plus difficile et plus honnête : à ces
tailles, la région que l'on sait retrancher -- celle qui entoure les points
déjà connus -- ne représente qu'une part négligeable de ce qui sépare
$\Rcal$ de $\Ecal$. Aucune quantité de coupes de ce genre n'y changera rien,
puisque aucune ne porte là où il le faudrait. Refermer cet écart demande un
relâché construit autrement. Nous ne savons pas le construire.

La seconde est la clôture elle-même : nos mesures n'ont jamais rencontré
deux points efficaces de même vecteur critères, ce qui suggère qu'une
condition structurelle -- vérifiable une fois pour toutes sur l'instance
plutôt qu'une fois par point d'archive -- rendrait la coupe entièrement
gratuite, le théorème~\ref{th:hauteur} n'en étant qu'une approximation
bon marché.

La troisième porte sur le raccourci de l'étape~1 de Zerdani \& Moulaï,
mis au jour par la section~\ref{sec:matchzm} : lorsque $\varphi$ est
linéaire, un seul programme entier donne son maximum sur $\Scal$, et si ce
point est efficace tout est fini. Notre $f$ étant fractionnaire, nous n'en
disposons pas -- le maximum de $f$ sur $\Scal$ demande un Dinkelbach
complet. Reste à savoir si un test bon marché, appliqué au \emph{premier}
itéré de ce Dinkelbach, en récupérerait une partie du bénéfice.

\end{document}
